\section{The method: SMC-BVS v3.1}\label{sec:method}

This section states the method exactly as specified (method v3.1 = METHODOLOGY\_v3 plus the MEET-v31 diff) and as
implemented in the frozen production code used for every study fit (code v2.2.1). Version numbers of the method
(v3, v3.1) and of the code (v2.1, v2.2, v2.2.1) are different series
\src{README.md:5-9}. The authoritative method text is \texttt{METHOD\_v3.1\_FINAL.md}, written as a diff against
\texttt{METHODOLOGY\_v3.md} (``M'' below) \src{NBH/branch_snapshots/MEET-v31/METHOD_v3.1_FINAL.md:1-9}. Definition IDs
\rid{DF01}--\rid{DF25} are anchored in Appendix~\ref{app:notation}; method-meeting decisions \rid{MD01}--\rid{MD13} are
defined in \S\ref{sec:method-md}.

\begin{plain}[the method in one paragraph]
Given $X$ ($n\times p$) and $y$, SMC-BVS computes the posterior over subsets $S$ of predictors under a unit-information
$g$-prior with a unit ridge, a vague inverse-gamma noise prior and a beta-binomial BB$(1,p)$ size prior capped at
$\lfloor n/2\rfloor$ variables. When the capped model space is small it enumerates it exactly. Otherwise it starts from an
exactly sampleable product-of-blocks approximation $\pi_0$ built from small correlation blocks, each enumerated exactly,
and tempers from $\pi_0$ to the posterior with an SMC bridge; a calibration phase builds the tempering schedule and a
fresh final run with frozen settings produces the reported HPM, PIPs and evidence \src{Docs/Methodology/METHODOLOGY_v3.md:25-33}.
v3.1 is ``v3 unchanged as a statistical method, plus a list of additions'' (parallel layer, optional stacking, LS0,
reporting changes, wording fixes) \src{NBH/branch_snapshots/MEET-v31/METHOD_v3.1_FINAL.md:26-41}.
\end{plain}

\begin{uncertain}[approval state of the method text]
\texttt{METHOD\_v3.1\_FINAL.md} still says: ``\textbf{Status: FINAL from MEET-v31 (chair), for HQ approval.} Not yet approved
by HQ or Felipe.'' \src{NBH/branch_snapshots/MEET-v31/METHOD_v3.1_FINAL.md:5}. The later plain summary says ``Approved by HQ
after the method meeting (MEET-v31), two independent reviews (Sol 6.1: a proof check and an audit), and the N9 kernel
check.'' \src{NBH/branch_snapshots/MEET-v31/FOR_FELIPE_FINAL_METHOD.md:3}; HQ wrote ``\textbf{The final method (v3.1) is
approved.}'' \pos{1584}. Felipe approved the four method-process decisions \pos{1304}
\src{NBH/START_HERE.md:16}; no source records Felipe approving the v3.1 method text itself. Treat v3.1 as HQ-approved;
the header line of the method file is stale.
\end{uncertain}

\subsection{Inputs, preprocessing, statuses}\label{sec:method-pre}

\textbf{Input} (M \S2): numeric $X\in\mathbb R^{n\times p}$, $y\in\mathbb R^n$, an integer seed $s\in[1,2^{31}-1]$, and the
job's resource allocation (memory, wall time; default 8 GiB) \src{Docs/Methodology/METHODOLOGY_v3.md:37-38}. v3.1 rewrites
Requirement (1) as: ``Runnable from X, y and one integer seed, with no statistical tuning beyond published defaults. The
job's allocation (memory, wall time) is a declared (c) input that decides delivery only. Every estimate on 𝒪 is a
function of (X, y, seed, binary).'' (item e8) \src{NBH/branch_snapshots/MEET-v31/METHOD_v3.1_FINAL.md:179}.

\textbf{Preprocessing} (\rid{DF04}; M \S4) \src{Docs/Methodology/METHODOLOGY_v3.md:80-93}:
\begin{enumerate}[label=(\arabic*)]
\item stop with \texttt{invalid\_input} if any entry is not finite, or $n<3$, or $p<1$ (v3; v3.1 item e3 proposes to admit
  $n=2$, see below);
\item drop column $j$ only if all its entries are \emph{exactly} equal; $p:=$ number kept; if $p=0$ return $\hat S=\varnothing$
  with status \texttt{no\_predictors};
\item if all entries of $y$ are exactly equal return $\hat S=\varnothing$, PIPs and evidence unavailable, status
  \texttt{y\_constant};
\item standardise stably: divide each column by its maximum absolute value, compute the mean, centre, divide by the sample SD
  (denominator $n-1$), so $\|x_j\|^2=n-1$; same for $y$, recording $s_y^2=\|Cy\|^2/(n-1)$, $C=I-\mathbf 1\mathbf 1^\top/n$;
\item report connected components of $\{|r_{jl}|\ge 1-10^{-10}\}$ (duplicates are kept);
\item resource preflight (M \S8.6).
\end{enumerate}
These rules make the target invariant to affine maps of $y$ ($a\ne0$) and of each column, and to column permutation up to
index-based tie-breaks \src{Docs/Methodology/METHODOLOGY_v3.md:92-93}. Code: \src{Scripts/Campaigns/2026-10-28_New_Campaign/code/core/v2_2_1/core/v3_preprocess.R:97-132}
(rejects $n<3$ at line 106; the near-constant ``anchored two-pass'' route at $\sqrt{\varepsilon_{\mathrm{mach}}}$ is at line 47).

\textbf{Statuses} (exactly one): \texttt{exact}, \texttt{completed}, \texttt{schedule\_limit}, \texttt{numerical\_failure},
\texttt{resource\_limit}, \texttt{y\_constant}, \texttt{no\_predictors}, \texttt{invalid\_input}
\src{Docs/Methodology/METHODOLOGY_v3.md:48-52}; v3.1 adds \texttt{degenerate\_n} (e3) and \texttt{unsupported\_size} (e7)
\src{NBH/branch_snapshots/MEET-v31/METHOD_v3.1_FINAL.md:80}. Item e3: ``Reject n = 0. n = 1 → status
\texttt{degenerate\_n} [\dots]. \textbf{n = 2 is admitted}: K = 1, B = 1, E = p + 1, so the exact branch runs.'' The $n=2$
output statement is ``seat-derived, unchecked'', and ``current preprocessing still rejects n < 3 and raw p = 0''
\src{NBH/branch_snapshots/MEET-v31/METHOD_v3.1_FINAL.md:182}.

\textbf{Flags}: \texttt{escalated}, \texttt{checks\_passed\_calibration}, \texttt{checks\_passed\_final},
\texttt{forced\_pilot\_steps}, \texttt{grid\_unconverged}, \texttt{local\_limit}, \texttt{low\_empirical\_ESS} (any final stage
below $0.5N$), \texttt{near\_tie\_truncated}, \texttt{local\_certified}, \texttt{global\_certified}
\src{Docs/Methodology/METHODOLOGY_v3.md:54-58}; v3.1 adds \texttt{ls0\_improved} and documents \texttt{global\_certified} as
``a heuristic numerical flag'' (e5) \src{NBH/branch_snapshots/MEET-v31/METHOD_v3.1_FINAL.md:81}
\src{NBH/branch_snapshots/MEET-v31/METHOD_v3.1_FINAL.md:184}.

\subsection{Target posterior: likelihood, priors, score, evidence}\label{sec:method-target}

Definitions \rid{DF01} (model, likelihood, $g$-prior, $\sigma^2$ prior), \rid{DF02} (size prior, cap), \rid{DF03} (score,
evidence, posterior). v3.1 keeps the target ``\textbf{unchanged} (Felipe's model)''
\src{NBH/branch_snapshots/MEET-v31/METHOD_v3.1_FINAL.md:84}.

\begin{stmt}{Docs/Methodology/METHODOLOGY_v3.md:97-108}
On the standardised data $(X,y)$, for $S\subseteq\{1,\dots,p\}$ with $k=|S|\le K$, $G_S=X_S^\top X_S$, $z_S=X_S^\top y$:
\[
y\mid\beta_S,\sigma^2,S\sim N_n(X_S\beta_S,\sigma^2I_n)\qquad\text{(no intercept; $y$ is centred)}
\]
\[
\beta_S\mid\sigma^2,S\sim N_k\bigl(0,\,n\sigma^2(G_S+I_k)^{-1}\bigr)\qquad\text{$g=n$ (unit information), $\tau=1$ (one pseudo-observation)}
\]
\[
\sigma^2\sim\mathrm{IG}(a_0,b_0),\quad a_0=b_0=0.01
\]
\[
\varpi(S)=w(k)\,1\{k\le K\}/W_K,\quad w(k)=B(1+k,2p-k)/B(1,p)=\tfrac12\big/\tbinom{2p-1}{k},\quad W_K=\textstyle\sum_{j\le K}\binom pj w(j)
\]
\[
H_S=(n+1)G_S+I_k,\qquad Q_S=y^\top y-n\,z_S^\top H_S^{-1}z_S,
\]
\[
\log m(S)=\log C_n-\tfrac12\bigl[\log\det H_S-\log\det(G_S+I_k)\bigr]-(a_0+n/2)\log(2b_0+Q_S),
\]
\[
C_n=(2b_0)^{a_0}\,\Gamma(a_0+n/2)\big/\bigl(\pi^{n/2}\,\Gamma(a_0)\bigr),\qquad
\pi(S\mid y)=\varpi(S)\,m(S)/Z,\qquad Z=\textstyle\sum_{|S|\le K}\varpi(S)\,m(S).
\]
\end{stmt}

The cap is $K=\min(p,\lfloor n/2\rfloor)$ and the model space is $\Omega=\{S:|S|\le K\}$
\src{Docs/Methodology/METHODOLOGY_v3.md:65} \src{NBH/branch_snapshots/MEET-THM/THEOREM_LIST_FINAL.md:64}. The shape is
$\nu=a_0+n/2$ and the theory cards write $\bar Q=2b_0+y^\top y$ \src{NBH/notation_and_conventions.md:142}.

\textbf{The score $F$.} Throughout the method and the theory, the \emph{score} of a model is the log unnormalised posterior
\[
F(S)=\log\varpi(S)+\log m(S),\qquad e^{F(S)}=\varpi(S)\,m(S),\qquad Z=\textstyle\sum_{S\in\Omega}e^{F(S)},
\]
(``the strict maximum of log ϖ(S) + log m(S)'' \src{Docs/Methodology/METHODOLOGY_v3.md:263}; ``$Z_{\mathcal T}=\sum_{S\in\mathcal T}e^{F(S)}=\sum_{\mathcal T}\varpi(S)\,m(S)$'' \src{NBH/notation_and_conventions.md:91}; code
\texttt{F = log varpi(|S|) + log m(S)} \src{Scripts/Campaigns/2026-10-28_New_Campaign/code/core/v2_2_1/core/cpp/v3_engine.h:89-98}).
So $F$ includes the prior; the marginal likelihood alone is $m(S)$. The Bayes factor against the empty model is
$\mathrm{BF}(S)=m(S)/m(\varnothing)$ \src{Docs/Methodology/METHODOLOGY_v3.md:176}. The final theory cards write
$m(U)=c_n[\det H_U/\det(G_U+I)]^{-1/2}(.02+Q_U)^{-\nu}$ and note ``The common positive $c_n$ cancels from posterior
ratios'' \src{NBH/branch_snapshots/MEET-THM/THEOREM_LIST_FINAL.md:65}.

\textbf{Reasons recorded for the target} (accepted by Felipe): $g=n$ is unit information and $\tau=1$ keeps every $m(S)$
finite for collinear $X_S$; IG$(0.01,0.01)$ is ``proper, as the project proofs require''; BB$(1,p)$ ``depends on p only;
prior mean size p/(p+1) ≈ 1; each extra variable costs log((2p−k−1)/(k+1)) nats''; the cap ``keeps near-interpolating
models out of the space; the ½ is a declared restriction, not needed for propriety''
\src{Docs/Methodology/METHODOLOGY_v3.md:111-119}. The v3 claim that the IG effect is negligible was narrowed in v3.1: its
pairwise effect ``relative to the gap [\dots] is not uniformly negligible [\dots]; negligibility needs residual-scale
assumptions; not negligible at tiny n'' \src{NBH/branch_snapshots/MEET-v31/METHOD_v3.1_FINAL.md:299}.

\textbf{Reduction statement} (consensus of the v3 consultants; ``to be proved, V1''): with $\theta=s_y\beta$ and
$v=s_y^2\sigma^2$, the posterior is that of the contrast regression of $U^\top y^r$ on $U^\top X$ with $n-1$ Gaussian errors,
slab $\theta_S\mid v,S\sim N(0,nv(G_S+I)^{-1})$ and $v\sim\mathrm{IG}(0.51,0.01s_y^2)$; ``This is not a claim that the
working evidence equals a raw-data marginal likelihood'' \src{Docs/Methodology/METHODOLOGY_v3.md:121-127}.

\textbf{Code check (prior/score).} The scorer sets $K=\min(p,\lfloor n/2\rfloor)$, $a_0=b_0=0.01$, $\nu=a_0+n/2$,
$\log C_n$, $\log w(k)=-\log 2-\log\binom{2p-1}{k}$, $\log W_K$ by log-sum-exp
\src{Scripts/Campaigns/2026-10-28_New_Campaign/code/core/v2_2_1/core/cpp/v3_score.h:39-72}, forms $H=(n+1)G_S+I$ and
$A=G_S+I$ by Cholesky and computes $Q$ by subtraction with recomputation by the non-negative form below
$10^{-8}y^\top y$ \src{Scripts/Campaigns/2026-10-28_New_Campaign/code/core/v2_2_1/core/cpp/v3_score.h:175-224}; Cholesky,
refactor, augmented QR, then \texttt{NumericalFailure}
\src{Scripts/Campaigns/2026-10-28_New_Campaign/code/core/v2_2_1/core/cpp/v3_score.h:227-238}. This matches M \S5 and \S6.13
\src{Docs/Methodology/METHODOLOGY_v3.md:286-289}. No divergence found.

\subsection{Budgets and the exact branch}\label{sec:method-budgets}

\begin{center}\small
\begin{tabular}{@{}lll@{}}
\toprule
Symbol & Definition (M \S3) & study ($n=200$; $p=500$ / $1000$)\\
\midrule
$K$ & size cap $\min(p,\lfloor n/2\rfloor)$ & 100 / 100\\
$B$ & block size $\min(\lfloor\sqrt n\rfloor,\lfloor\sqrt p\rfloor,B^*)$, $B^*=16$ & 14 / 14\\
$\ell$ & $\lceil\log_2(n+p+1)\rceil$ (local-search and rebuild budget) & 10 / 11\\
$L$ & $\max(4,\lceil\log_2(n+1)\rceil)$ (burn-in rounds at $\varphi=1$) & 8 / 8\\
$N_0$ & $\max(2500,n,p)$ final particles per island & 2500 / 2500\\
$H_0$ & $\lceil 512N_0/2500\rceil$ calibration particles (pilot and each diagnostic island) & 512 / 512\\
$T_0$ & $\max(1000,\ell^2)$ stage ceiling & 1000 / 1000\\
$R$ & 4 islands (final and diagnostic) & 4 / 4\\
tier $e$ & $e\in\{0,1\}$; $(N_e,H_e,T_e)=4^e(N_0,H_0,T_0)$ & \\
\bottomrule
\end{tabular}
\end{center}
Source \src{Docs/Methodology/METHODOLOGY_v3.md:62-76}; ``These are \textbf{declared resource budgets}, not accuracy
guarantees.'' v3.1 keeps the values and relabels classes: $B^*=16$ is ``a declared \textbf{(c) budget for exhaustive
tables}'' (e1) \src{NBH/branch_snapshots/MEET-v31/METHOD_v3.1_FINAL.md:180}; ``the ℓ² term exceeds 1,000 exactly when
n + p ≥ 2³¹'' (w4) \src{NBH/branch_snapshots/MEET-v31/METHOD_v3.1_FINAL.md:195}. Code:
\src{Scripts/Campaigns/2026-10-28_New_Campaign/code/core/v2_2_1/core/v3_preprocess.R:19-33} (exact integer
$\lfloor\sqrt{\cdot}\rfloor$ and $\lceil\log_2\cdot\rceil$; matches).

\textbf{Exact branch} (M \S6.2): if $E=\sum_{k\le K}\binom pk\le(5H_0+4N_0)\,p$ (saturating integer arithmetic), score every
model with $|S|\le K$ in a streaming pass (running best in the strict HPM order; log-sum-exp for $Z$ and PIP numerators);
status \texttt{exact} \src{Docs/Methodology/METHODOLOGY_v3.md:156-160}. ``Because the threshold grows with N₀ ≥ n, every
fixed-p problem is eventually enumerated as n grows.'' v3.1 corrects ``all p ≤ 3,000 at n ≤ 5'' to ``all p at n ≤ 5'' (w2)
\src{NBH/branch_snapshots/MEET-v31/METHOD_v3.1_FINAL.md:193}. Code: 128-bit saturating count
\src{Scripts/Campaigns/2026-10-28_New_Campaign/code/core/v2_2_1/core/cpp/v3_core.cpp:97-110}; enumeration
\src{Scripts/Campaigns/2026-10-28_New_Campaign/code/core/v2_2_1/core/cpp/v3_core.cpp:526-570}. On this branch
\texttt{global\_certified} is set iff $F(\mathrm{HPM})-\eta(\mathrm{HPM})>\max_{S\ne\mathrm{HPM}}[F(S)+\eta(S)]$ with the
scorer's heuristic error expression $\eta$ \src{Scripts/Campaigns/2026-10-28_New_Campaign/code/core/v2_2_1/core/v3_exact.R:3-17};
v3.1 (e5) labels this ``a \textbf{heuristic numerical margin}'' \src{NBH/branch_snapshots/MEET-v31/METHOD_v3.1_FINAL.md:184}.
The 2,400-fit study used the SMC branch, where the flag is always FALSE
\src{Scripts/Campaigns/2026-10-28_New_Campaign/code/core/v2_2_1/core/v3_run.R:147}.

\subsection{The block start \texorpdfstring{$\pi_0$}{pi0}}\label{sec:method-start}

Definitions \rid{DF05} (blocks and tables), \rid{DF06} (start). v3.1 keeps blocks, tables and start ``\textbf{unchanged}'';
the table stage runs in parallel (\#4) \src{NBH/branch_snapshots/MEET-v31/METHOD_v3.1_FINAL.md:87}.

\begin{stmt}{Docs/Methodology/METHODOLOGY_v3.md:164-177}
\begin{itemize}
\item \textbf{Clustering:} complete linkage on the distance $1-|r_{jl}|$ computed from $X$; ties between candidate merges broken
  lexicographically on the sorted original-index lists.
\item \textbf{Cut:} top-down from the root; any node with more than $B$ members is split into its two children (visited smallest
  index first) until every retained cluster has at most $B$ members.
\item \textbf{Packing:} next-fit in dendrogram leaf order, whole clusters only (a retained cluster is never split), into blocks
  $C_1,\dots,C_J$ of at most $B$ members. $J\ge\lceil p/B\rceil$ (packing leaves gaps).
\item \textbf{Tables:} for each block, every sub-model $U\subseteq C_b$ is scored against the full $y$; $r_b(U)=m(U)/m(\varnothing)$
  (table exponent 1). Since $B\le K$ for $n\ge3$, no sub-model exceeds the cap. Every table score counts for the HPM. Table
  scores are not inserted in the score cache.
\item \textbf{Start:} $\gamma_0(S)=\varpi(S)\prod_b r_b(S\cap C_b)$, $\pi_0=\gamma_0/\bar Z_0$. $\bar Z_0$ is computed exactly by
  convolving the block size profiles $e_b(j)=\sum_{|U|=j}r_b(U)$ up to $K$ with the size prior (valid because $\varpi$ depends on
  $S$ only through $|S|$). Exact iid draws: draw the total size, then block sizes backwards, then each block's sub-model within
  its size class.
\item \textbf{Discrepancy:} $d(S)=\log\mathrm{BF}(S)-\sum_b\log r_b(S\cap C_b)$, $\mathrm{BF}(S)=m(S)/m(\varnothing)$; bridge
  $\gamma_\varphi(S)=\gamma_0(S)e^{\varphi d(S)}$, $\varphi\in[0,1]$; $\gamma_1\propto$ the posterior; $Z=m(\varnothing)\sum_S\gamma_1(S)$.
\end{itemize}
\end{stmt}

\textbf{Block-size rule and no correlation cut.} $B=\min(\lfloor\sqrt n\rfloor,\lfloor\sqrt p\rfloor,B^*)$ is ``Felipe's rule;
B ≤ K and B ≤ p'' \src{NBH/branch_snapshots/MEET-v31/METHOD_v3.1_FINAL.md:307}. w5: ``The dendrogram is cut by node size
only (≤ B). There is no correlation-height threshold.'' \src{NBH/branch_snapshots/MEET-v31/METHOD_v3.1_FINAL.md:196}. e6
records the study's actual block counts: ``J = 46–60 (p = 500) and 95–120 (p = 1,000) versus proxies 36 and 72''
\src{NBH/branch_snapshots/MEET-v31/METHOD_v3.1_FINAL.md:185}. The theory's partition hypothesis \textbf{SIZE}
$g_{\max}\le B<2g_{\min}$ is ``a hypothesis of one partition theorem [\dots], not a method parameter. The study's g = 5 is the
theorem obligation PART-pairing'' \src{NBH/branch_snapshots/MEET-v31/METHOD_v3.1_FINAL.md:180} (obligation \rid{OB08};
card \rid{TH11}).

\textbf{Exact sampling (\texttt{draw\_pi0}) and $\bar Z_0$ in code.} Tables score every $U\subseteq C_b$ with
$|U|\le K$ against the full $y$ \src{Scripts/Campaigns/2026-10-28_New_Campaign/code/core/v2_2_1/core/cpp/v3_start.h:73-98};
size profiles $\log e_b(j)$, partial convolutions $\log F_b(k)$ truncated at $K$, $\log\bar Z_0=\log\sum_k\varpi(k)F_J(k)$, and
the backward conditional $P(k_b=j\mid r)=e_b(j)F_{b-1}(r-j)/F_b(r)$
\src{Scripts/Campaigns/2026-10-28_New_Campaign/code/core/v2_2_1/core/cpp/v3_start.h:100-151}; the draw uses one uniform for
the total size, one per block $b=J,\dots,1$ while the remainder is positive, and one per non-empty block for the
sub-model \src{Scripts/Campaigns/2026-10-28_New_Campaign/code/core/v2_2_1/core/cpp/v3_start.h:153-176}. Blocks:
\src{Scripts/Campaigns/2026-10-28_New_Campaign/code/core/v2_2_1/core/cpp/v3_blocks.h:39-144} (complete linkage with the
union-lexicographic tie rule; left child = child with smaller minimum index; next-fit packing). All match M \S6.3.
Exactness of \texttt{draw\_pi0} and $\bar Z_0$ is part of N9 (\rid{OB09}), PASSED \src{NBH/branch_snapshots/MEET-v31/MINUTES.md:150}
(card \rid{TH12}; clause \rid{TH01.a}).

\textbf{Start mode $\tilde S_0$.} $\tilde S_0=\arg\max\gamma_0$, ``computed exactly by the block DP, with the size-then-lexicographic
tie rule'' \src{NBH/branch_snapshots/MEET-v31/METHOD_v3.1_FINAL.md:150}. The header function \texttt{start\_mode} computes
per-block per-size maxima, a max-plus convolution up to $K$, picks the size by strict improvement (``ties keep the smaller
size'') and enumerates all optimal choices keeping the lexicographically smallest set, capped at $10^6$ sets with a
truncation flag \src{Scripts/Campaigns/2026-10-28_New_Campaign/code/core/v2_2_1/core/cpp/v3_start.h:178-240}. e4: ``On
truncation, use a deterministic tie-aware DP that returns the prescribed smallest-size, lexicographically smallest
maximiser without enumerating ties'' \src{NBH/branch_snapshots/MEET-v31/METHOD_v3.1_FINAL.md:150}; ``The invariant `DP output
= the prescribed choice' is an OPEN obligation (N3)'' \src{NBH/branch_snapshots/MEET-v31/METHOD_v3.1_FINAL.md:183}
(\rid{OB03}). In the study, ``0/2,400 tie truncations; all 2,400 had a single tie''
\src{NBH/branch_snapshots/MEET-v31/METHOD_v3.1_FINAL.md:150}.

\textbf{Dilution (cite, do not derive).} The tables score each block's sub-models against the \emph{full} $y$ with the
noise variance integrated out per sub-model, so a block's table does not see the other blocks' signal except as noise.
O-DILUTE: ``Seen from a truth, the v3 table's own-σ² integration \emph{dilutes}; seen from a null, it \emph{calibrates} the
leakage'' \src{NBH/branch_snapshots/O-DILUTE/THEORY.md:114-115}; mechanism ``σ²-dilution of the truths' table log-odds''
with ``About 7 of the 14 truths are invisible at n = 200 for every c''
\src{NBH/03_results_and_conclusions.md:269}. The exact identity $-\log\pi_0(S_0)=V+C-\log A(|S_0|)$ is card M3(d)
\src{NBH/branch_snapshots/MEET-THM/THEOREM_LIST_FINAL.md:206} (\rid{TH03}). The common plug-in $\sigma^2$ alternative was
rejected (it ``floods the start'') \src{NBH/03b_reasoning_and_traces.md:83-88}.

\textbf{Item 1 (start ablation).} The start is kept under Outcome 2: ``the start is presented as the exactly computable,
theorem-bearing local step (S1/S2) plus the parallel table stage, with \textbf{no depth or speed claim}''
\src{NBH/branch_snapshots/MEET-v31/METHOD_v3.1_FINAL.md:118} (\rid{MD01}).

\subsection{Bridge, neighbour graph and move kernel}\label{sec:method-kernel}

Definitions \rid{DF07} (discrepancy and bridge), \rid{DF08} (moves and round). The bridge family is
$\gamma_\varphi=\gamma_0e^{\varphi d}$ (above); the theory cards write the normalised family $q_t\propto\gamma_0e^{td}$
\src{NBH/branch_snapshots/MEET-THM/THEOREM_LIST_FINAL.md:67} and the registrar's canonical notation is
$q_\varphi:=\gamma_\varphi/\bar Z_\varphi$ \src{NBH/notation_and_conventions.md:55}.

\textbf{Neighbour graph} (M \S6.4): $r_n=\sqrt{2\log p/n}$; for each $j$ the partners $l\ne j$ with $|r_{jl}|>r_n$, always
including $j$'s single largest-$|r|$ partner, keeping at most the $B$ largest (ties by original index); then the symmetric
union $N(j)$ (degrees may exceed $B$); no edges when $p=1$ \src{Docs/Methodology/METHODOLOGY_v3.md:181-183}. w3: when
$r_n\ge1$ (i.e.\ $p>e^{n/2}$) ``the graph is the symmetric 1-nearest-neighbour graph''
\src{NBH/branch_snapshots/MEET-v31/METHOD_v3.1_FINAL.md:194}. Code matches
\src{Scripts/Campaigns/2026-10-28_New_Campaign/code/core/v2_2_1/core/cpp/v3_blocks.h:146-173}.

\begin{stmt}{Docs/Methodology/METHODOLOGY_v3.md:187-201}
A round is the list of $m=p+B+2$ labelled elementary steps --- the $p$ Gibbs coordinate updates, one uniform swap, $B$ neighbour
swaps, one independence proposal --- applied \textbf{in a uniformly random order drawn independently of the particle's state,
afresh for every particle and every round}.
\begin{itemize}
\item \textbf{Gibbs coordinate $j$:} draw $j$'s inclusion from its conditional under $\gamma_\varphi$, with an addition at $|S|=K$
  excluded. Record the conditional inclusion probability $o_j$ immediately before the draw (used for PIPs in measured rounds).
\item \textbf{Uniform swap:} skipped (self-loop) if $|S|\in\{0,p\}$; delete $j$ uniformly from $S$, add $l$ uniformly from outside;
  accept with $1\wedge\gamma_\varphi(S')/\gamma_\varphi(S)$.
\item \textbf{Neighbour swap:} skipped if $S=\varnothing$; delete $j$ uniformly from $S$; if $N(j)\setminus S=\varnothing$ reject;
  else add $l$ uniformly from $N(j)\setminus S$; accept with
  $1\wedge[\gamma_\varphi(S')/\gamma_\varphi(S)]\cdot|N(j)\setminus S|/|N(l)\setminus S'|$.
\item \textbf{Independence proposal:} $S'\sim\pi_0$ (exact draw), accepted with $1\wedge e^{\varphi[d(S')-d(S)]}$. Every proposal is
  scored, accepted or not.
\end{itemize}
Each step is reversible with respect to $\gamma_\varphi$; the composition in a random order whose law is invariant under reversal is
self-adjoint (checked by all three consultants; V5). An individual realised order is not.
\end{stmt}

\textbf{Invariance status.} N9 (``kernel reversibility and self-adjointness, exact start sampler'') was a closure gate of
the method meeting; ``\textbf{N9 PASSED} --- all seven sub-claims PROVED (candidate, author-only), with exact numerical
verification on the production engine headers and failing negative controls''
\src{NBH/branch_snapshots/MEET-v31/MINUTES.md:150}. The sub-claim count is UNCERTAIN: other records give 7 \pos{1550} and 8 \pos{1593} (see the uncertain boxes under \rid{TH12} and \rid{OB09}). See \rid{OB09}, \rid{TH12} (S-KERNEL) and \rid{TH01.a} for the statement
and its exact status words. v3.1 keeps the round unchanged; ``Changing it is WORKHORSE''
\src{NBH/branch_snapshots/MEET-v31/METHOD_v3.1_FINAL.md:89}.

\textbf{Code check (kernels).} \src{Scripts/Campaigns/2026-10-28_New_Campaign/code/core/v2_2_1/core/cpp/v3_engine.h:111-215}:
Gibbs computes $o_j=\mathrm{plogis}(\log\gamma_\varphi(S\cup j)-\log\gamma_\varphi(S\setminus j))$, sets $o_j=0$ for an infeasible
addition at $|S|=K$, records $o_j$ before the draw, and counts both configurations as scored (``SR'');
uniform swap self-loops at $|S|\in\{0,p\}$; neighbour swap self-loops at $S=\varnothing$, rejects without scoring when
$N(j)\setminus S=\varnothing$, and includes the Hastings factor $|N(j)\setminus S|/|N(l)\setminus S'|$; the independence step
draws $S'\sim\pi_0$ and accepts with $\exp\{\varphi(d(S')-d(S))\}\wedge1$; the order is a Fisher--Yates permutation of the $m$
labels drawn afresh per particle and round. All match M \S6.5. The exact uniform-consumption order is documented in
\src{Scripts/Campaigns/2026-10-28_New_Campaign/code/core/v2_2_1/core/v3_run.R:24-40}.

\subsection{Resampling, pilot schedule, candidate grid and calibration}\label{sec:method-calib}

Definitions \rid{DF09} (resampling, pilot, ESS 0.90), \rid{DF10} (candidate set, grid, rebuild loop), \rid{DF11}
(calibration record, tiers, final runs).

\begin{stmt}{Docs/Methodology/METHODOLOGY_v3.md:205-227}
\begin{itemize}
\item \textbf{Resampling:} fresh multinomial resampling after every reweighting, including the terminal stage ($O(N)$ per stage with
  a linear-time categorical sampler).
\item \textbf{Pilot:} $H_e$ particles from $\pi_0$; at each stage the next temperature is the largest $\varphi'$ with population ESS of
  the weights $e^{(\varphi'-\varphi)d}$ at least $0.90\cdot H_e$ (weight CV$^2\le1/9$), found by 60 bisection halvings; exploratory floor
  $1/T_e$ (forced steps counted as \texttt{forced\_pilot\_steps}); ends exactly at $\varphi=1$; then $L$ burn-in rounds, 4 measured
  rounds and a local search.
\item \textbf{V:} the union of the top $20p$ scored models by $\log\gamma_0$, the top $20p$ scored models by $\log\gamma_1$ (the target),
  and $\varnothing$; deduplicated; list ties lexicographic; frozen during each construction.
\item \textbf{Restricted log-normaliser:} $F_V(t)=\log\sum_{S\in V}\gamma_0(S)e^{t\,d(S)}$.
\item \textbf{Restricted ESS} between temperatures $u<v$: $\rho_V(u,v)=\exp\{2F_V(v)-F_V(u)-F_V(2v-u)\}$, the population ESS fraction
  of the weights $e^{(v-u)d}$ under the restricted law $q_u(\cdot\mid V)$ (evaluating $F_V$ beyond 1 is a finite-sum moment, not a
  production temperature).
\item \textbf{Construction:} start at $u=0$ with the pilot temperatures as targets. For the next pilot temperature $b$: if
  $\rho_V(u,b)\ge0.90$ append $b$; otherwise bisect on $(u,b]$ with $u$ fixed, at most 60 halvings, and append the largest tested
  endpoint that passes; continue from the appended point. Validate strict numerical progress and the inequality at every
  appended point. Non-progress, failed validation, or more than $T_e$ intervals $\to$ the construction fails.
\item \textbf{Rebuild loop:} after each set of diagnostic islands, update $V$ and rebuild; stop when the rebuilt grid equals the
  previous one; at most $\ell$ iterations [\dots]. If not converged: flag \texttt{grid\_unconverged} (a calibration trigger); at
  tier 1 the last grid is used and the flag is reported. If a construction fails: trigger escalation at tier 0; status
  \texttt{schedule\_limit} at tier 1.
\end{itemize}
\end{stmt}
[\dots] = the lead's ruling that convergence is judged on the grid, not on $V$, with Fable-T3's dissent (judge on $V$)
\src{Docs/Methodology/METHODOLOGY_v3.md:224-225}. Purpose and limits: the grid inserts temperatures ``exactly where models
already scored would make the bridge jump''; it ``controls restricted one-step weight dispersion only; it does not control
full-space ESS, particle occupancy, mixing, or modes never scored'' \src{Docs/Methodology/METHODOLOGY_v3.md:228-236}.

\textbf{Calibration and escalation} (M \S6.8--6.9): calibration at tier $e$ runs the pilot and, per rebuild iteration, 4
diagnostic islands of $H_e$ particles on the current grid, each followed by $L$ burn-in rounds, 4 measured rounds and a local
search; streams 0 and 1--4 (tier 0), 5 and 6--9 (tier 1); score leaders are kept across tiers
\src{Docs/Methodology/METHODOLOGY_v3.md:240-242}. Triggers (any one $\Rightarrow$ at most one escalation to tier 1, decided before
any final randomness): E1 (some diagnostic island's best score more than 1 nat below the overall best), E2 (SD of the 4
diagnostic $\log\hat Z$ exceeds 1 nat), E3 (a local search improved the best), forced pilot steps, \texttt{grid\_unconverged},
a failed grid construction, \texttt{local\_limit}; after the final run E1--E3 and the empirical-ESS check are recomputed as
flags only \src{Docs/Methodology/METHODOLOGY_v3.md:249-254}. v3.1 keeps these ``as calibration triggers, which define the law of
𝓒'' and labels them heuristics (\rid{MD09}) \src{NBH/branch_snapshots/MEET-v31/METHOD_v3.1_FINAL.md:93}.

\begin{stmt}{NBH/branch_snapshots/MEET-v31/METHOD_v3.1_FINAL.md:170}
\textbf{Calibration contract.} ``The calibration output 𝓒 (grid, budgets $N_e$, tier, $R'$, $V$-independent settings) is frozen
before any final randomness. Validity of the final estimates (unbiased $\hat Z_r$, PIP identities) holds conditionally on
\textbf{any} frozen 𝓒, whatever happened in calibration. The grid satisfies Theorem G's cap. There are $\le\ell$ rebuilds per tier
($\le2\ell$ in total with escalation, \texttt{v3\_calib.R:85, 141--149}). The restricted ESS controls one-step weight dispersion on
$V$ only. The statuses \texttt{schedule\_limit}, \texttt{grid\_unconverged} and \texttt{escalated} are distinguished. Rebuild
convergence, $P(S_0\in V)$ and $P(\texttt{schedule\_limit})$ are \textbf{OPEN theorem obligations, not inherent limitations}.
Rebuilds, tier, $T$ and convergence are tabulated POST HOC from saved fits.''
\end{stmt}

\textbf{Code check (calibration).} Pilot next temperature (ESS $\ge0.90H$, 60 halvings, floor $1/T_e$, snap within
$10^{-12}$ of 1) and grid construction (targets = pilot temperatures, 60 halvings, validation, $T_e$ cap)
\src{Scripts/Campaigns/2026-10-28_New_Campaign/code/core/v2_2_1/core/cpp/v3_calib.h:25-98}; rebuild loop of at most $\ell$
iterations with convergence by grid equality \src{Scripts/Campaigns/2026-10-28_New_Campaign/code/core/v2_2_1/core/v3_calib.R:84-119};
at most one escalation, re-preflight before tier 1, freeze of (grid, tier, $N$, $L$, inherited leaders and HPM)
\src{Scripts/Campaigns/2026-10-28_New_Campaign/code/core/v2_2_1/core/v3_calib.R:135-175}. The frozen record in code also
carries the \emph{inherited} calibration stores (top store and both leader stores) and the calibration HPM with its
provenance, which the final run offers to its stores \src{Scripts/Campaigns/2026-10-28_New_Campaign/code/core/v2_2_1/core/cpp/v3_core.cpp:731-743}
\src{Scripts/Campaigns/2026-10-28_New_Campaign/code/core/v2_2_1/core/cpp/v3_core.cpp:443-466}. Matches M \S6.6--6.9.

\subsection{Final islands, streams, stacking and the parallel layer}\label{sec:method-islands}

Definitions \rid{DF11}, \rid{DF12} (islands and stacking). \textbf{Final run} (M \S6.8): 4 islands $\times N_e$ particles from
fresh iid $\pi_0$ draws, the frozen grid, multinomial resampling, one move round per stage, then $L$ burn-in rounds and 4
measured rounds at $\varphi=1$, then a local search; streams 10--13, independent of all calibration randomness; no forced
steps \src{Docs/Methodology/METHODOLOGY_v3.md:243-245}. ``Final observations never revise the schedule and never select
replacement estimates.'' \src{Docs/Methodology/METHODOLOGY_v3.md:151-152}. In code, each stream is a xoshiro256** generator
seeded by SplitMix64 from $x=(\mathrm{seed}\ll32)+\mathrm{stream}$; uniforms are $(\mathrm{next}\gg11)\,2^{-53}$
\src{Scripts/Campaigns/2026-10-28_New_Campaign/code/core/v2_2_1/core/v3_run.R:12-23}.

\begin{stmt}{NBH/branch_snapshots/MEET-v31/METHOD_v3.1_FINAL.md:142}
\textbf{Optional stacking $R'=4k$ at $N_e$.}
\begin{itemize}
\item ``$R'=4$ by default. Optionally $R'=4k$, at \textbf{unchanged $N_e$}, declared before final randomness (before stream 10).''
\item ``Final islands use streams $10,\dots,9+R'$. Islands 1--4 keep streams 10--13. \textbf{$N_e$ is never split} across more islands.''
\item ``With $R'>4$, an \textbf{extra local search runs from the islands-1--4 centre $c_4$}, besides the search from the expanded centre.
  The post-search scored set is then nested and HPM$(R')\succeq$ HPM(4) (Cor W-R with the extra search, PC1 claim 6).''
\item ``Two records are kept. The \textbf{expanded record} covers all $R'$ islands, pooled $\hat Z$ over $R'$ and the expanded store; it
  carries the run's reported flags. The \textbf{baseline record} covers islands 1--4, with its own baseline HPM, the search from $c_4$
  and its own E1--E3, low-ESS and \texttt{local\_certified}. It is kept whenever baseline reproducibility is claimed.''
\item ``Each record has its own certificate card (\#13).''
\end{itemize}
\end{stmt}
AMENDMENT\_HQ\_02 adds the optional configuration $R'_n=4\lceil\log_2(n+p+1)\rceil$ (default stays $R'=4$); the rate
$4\lceil(n+p)^6/4\rceil$ is ``a sufficient rate inside theorem M1 only, not a method configuration''
\src{NBH/branch_snapshots/MEET-v31/METHOD_v3.1_FINAL.md:425} \src{NBH/branch_snapshots/MEET-v31/AMENDMENT_HQ_02.md:9-12}.
An adaptive $R'$ is not allowed (``R′ must be declared in advance'') \src{NBH/branch_snapshots/METH-AUDIT/REPORT.md:214-215}.
Splitting $N_e$ into many islands (C8) was rejected \src{NBH/branch_snapshots/MEET-v31/METHOD_v3.1_FINAL.md:274}.

\textbf{Parallel execution layer} (\rid{MD04}) \src{NBH/branch_snapshots/MEET-v31/METHOD_v3.1_FINAL.md:131-136}:
(a) the $R'$ final islands run as concurrent \emph{process} workers with their own score cache and ledger, merged by
Lemma S before the final local search; ``Output is bitwise-identical to serial execution (Prop W). `Bitwise' is scoped to
one binary (\texttt{Makevars.v3}: \texttt{-ffp-contract=off -march=x86-64-v2}), single-thread BLAS and process workers'';
(b) the 4 diagnostic islands per rebuild run concurrently --- ```Bitwise' may \textbf{not} be claimed for (b) until W-cal
has been executed''; (c) local-search scan candidates may be scored in parallel with a max-reduction that keeps the
lexicographic tie rule; (h) block tables are scored in parallel, with no change of floating-point summation order under a
bitwise claim; workers are processes (threads would need private counters). (a) ``has been executed bitwise on smoke plus
3/3 study designs (CONC)'' \src{NBH/branch_snapshots/MEET-v31/METHOD_v3.1_FINAL.md:46}; the speed figure 1.9$\times$ ``is a 2-core
measurement'' \src{NBH/branch_snapshots/MEET-v31/METHOD_v3.1_FINAL.md:56}.

\textbf{GUARD-W} (\#6): ``Any worker stop, LS0 stop or guard stop ⇒ no output. The fit returns the corresponding unavailable
status with all estimate fields NA (\#15), never a partial estimate. [\dots] A completed LS0 that ends on its move cap is not
a stop: it returns output with \texttt{local\_limit}, and L1 = `n/a' (\#10). Lemma RO (a guard that reads no uniforms is
stop-only) stays as recorded (XF VERIFIED). That the v3.1 worker, merge and LS0 code satisfy GUARD-W is a
\textbf{code-audit obligation} (\S6), not a mathematical gap.'' \src{NBH/branch_snapshots/MEET-v31/METHOD_v3.1_FINAL.md:136}.

\textbf{Depth and cost} (\#19, \rid{MD06}): ``With island parallelism, the particle-round contribution to depth is
$N_e\cdot(T+L+4)\cdot(p+B+2)$ \textbf{sequential elementary-step slots} per island at tier e'' and ```Depth free of N, H, R'
is \textbf{not} claimed for v3.1''; the single-tier calibration/final trajectory ratio $(1+4I)H_0/(4N_0)=0.0512+0.2048\cdot I$ is
``\textbf{seat-derived, unchecked}'' \src{NBH/branch_snapshots/MEET-v31/METHOD_v3.1_FINAL.md:171}. Local-search cost bound:
scan size $p+k^*(p-k^*)$ with $k^*=\min(K,\lfloor p/2\rfloor)$, times $(\ell+1)$ scans
\src{NBH/branch_snapshots/MEET-v31/METHOD_v3.1_FINAL.md:99}. Declared v3.2 path (not shipped): C6 per-particle counter-based
streams, C7 Gram-free scorer, C11 faster exact linkage \src{NBH/branch_snapshots/MEET-v31/METHOD_v3.1_FINAL.md:259-265}.

\subsection{Outputs}\label{sec:method-outputs}

\subsubsection{Evidence \texorpdfstring{$\hat Z$}{Zhat}}
Definition \rid{DF13}.
\begin{stmt}{Docs/Methodology/METHODOLOGY_v3.md:279-282}
Final island $r$: $\hat Z_r=m(\varnothing)\,\bar Z_0\prod_t(1/N_e)\sum_i\exp\{(\varphi_t-\varphi_{t-1})\,d(S_i^{(t-1)})\}$, with $W_K$ inside
$\gamma_0$ exactly once. Pooled $\hat Z=(1/4)\sum_r\hat Z_r$ via log-sum-exp; 95\% lower bound $0.05\cdot\hat Z$ (Markov; one-sided, not an
accuracy interval). Reported on the full working scale and relative to $m(\varnothing)$. Given the frozen calibration, each $\hat Z_r$
is unbiased (existing project results P6, P7); $\log\hat Z$ is not.
\end{stmt}
v3.1 pools over the $R'$ islands (``pooled Ẑ over R′'') \src{NBH/branch_snapshots/MEET-v31/METHOD_v3.1_FINAL.md:142}, adds
$Z_{\mathcal T}=\sum_{S\in\mathcal T}e^{F(S)}$ with ``Ẑ < Z\_𝒯 implies Ẑ < Z'' (\#14)
\src{NBH/branch_snapshots/MEET-v31/METHOD_v3.1_FINAL.md:161}, and the card bound $\max\{Z_{\mathcal T},0.05\hat Z\}$ which ``keeps 95\%
coverage on the output event'' \src{NBH/branch_snapshots/MEET-v31/METHOD_v3.1_FINAL.md:173}. Paper sentence (\#17): ``The
probability that a run returns a card with a false claim is at most 5\%.'' \src{NBH/branch_snapshots/MEET-v31/METHOD_v3.1_FINAL.md:169}.
Unbiasedness given $\mathcal C$ is clause \rid{TH01.a}. Code: $\log\hat Z_r=\log m(\varnothing)+\log\bar Z_0+\sum_t$ log-increment,
increment $=\mathrm{LSE}(\Delta\varphi\,d_i)-\log N$, computed before resampling
\src{Scripts/Campaigns/2026-10-28_New_Campaign/code/core/v2_2_1/core/cpp/v3_engine.h:300-323}; pooled $\log\hat Z$, $\log(0.05)+\log\hat Z$,
and the value relative to $m(\varnothing)$ \src{Scripts/Campaigns/2026-10-28_New_Campaign/code/core/v2_2_1/core/v3_run.R:117-124}. Matches.

\subsubsection{PIPs and MPMs}
Definition \rid{DF14}. In each island, the RB-PIP averages the recorded Gibbs conditionals $o_j$ over particles and the 4
measured rounds \src{Docs/Methodology/METHODOLOGY_v3.md:271-272}.
\begin{stmt}{NBH/branch_snapshots/MEET-v31/METHOD_v3.1_FINAL.md:158}
\begin{itemize}
\item ``\textbf{Primary:} $\widehat{\mathrm{PIP}}_j=\sum_r\hat Z_r\widehat{\mathrm{PIP}}_{rj}/\sum_r\hat Z_r$ over the $R'$ final islands.''
\item ``\textbf{Secondary (named):} the equal-weight mean over islands.''
\item ``Report $\mathrm{MPM}_{\hat Z}=\{j:\text{primary }\widehat{\mathrm{PIP}}_j>\tfrac12\}$ and $\mathrm{MPM}_{\mathrm{eq}}$. Either may exceed $K$ and
  is reported as computed.''
\item ``Report $R_{\mathrm{eff}}=(\sum_rW_r)^2/\sum_rW_r^2$ with $W_r\propto\hat Z_r$.''
\item ``Neither normalised estimate is generally unbiased at a finite budget.''
\item ``The R-scaling guarantee is stated for the primary PIP; the equal-weight PIP keeps its single-island bias for every R (Prop P(i)).''
\end{itemize}
\end{stmt}
This \emph{swaps} v3's pooling, where the equal-weight mean was primary and the $\hat Z$-weighted one secondary
\src{Docs/Methodology/METHODOLOGY_v3.md:272-275} (\rid{MD08}). Code (v2.2.1) also computes \emph{plain} PIPs: per island the
fraction of particles containing $j$ at the end of each of the 4 measured rounds
\src{Scripts/Campaigns/2026-10-28_New_Campaign/code/core/v2_2_1/core/cpp/v3_engine.h:350-371}, pooled with equal weights as
\texttt{pip\_plain}, \texttt{mpm\_plain} \src{Scripts/Campaigns/2026-10-28_New_Campaign/code/core/v2_2_1/core/v3_run.R:167-169}.
These are the plain measured-endpoint averages $H_r(h)$ of \S\ref{sec:method-genh} for $h=1\{j\in S\}$, equal-weight pooled
(this identification is an editorial reading of the two definitions, not a source statement).

\subsubsection{HPM, local search, LS0}
Definitions \rid{DF15} (HPM and local search), \rid{DF16} (LS0).
\begin{stmt}{Docs/Methodology/METHODOLOGY_v3.md:258-267}
\begin{itemize}
\item \textbf{Local search:} from the current best scored model (the centre), scan every feasible flip and swap (additions beyond $K$
  excluded) with the centre frozen during the scan; move to the best strict improvement (lexicographic ties); at most $\ell$ moves,
  then one further scan; stop earlier when a scan finds no improvement. Every evaluated model updates the global best. If the
  last scan still improves, set \texttt{local\_limit} (no restart). \texttt{local\_certified} only when the returned model is a scanned
  centre with no improving neighbour. The search never alters particles.
\item \textbf{HPM:} the strict maximum of $\log\varpi(S)+\log m(S)$ over every model ever scored (tables, pilot, calibration, final rounds,
  rejected proposals, independence proposals, local searches), kept as a running best that survives cache eviction. Only exactly
  equal scores invoke the tie rule: smallest $|S|$, then lexicographically smallest index set. A near-tie list (models within
  $10^{-9}$ nat of the best, drawn from the stored top $40p+1$) is reported; it never reorders, and \texttt{near\_tie\_truncated} is set
  if the list may be incomplete.
\end{itemize}
\end{stmt}
Equivalently the HPM is the $\succ$-maximum of the scored set $\mathcal S$ for the strict total order
$S\succ S'\iff(F(S),-|S|,\cdot)$ is larger, the last component being ``lexicographically smaller index set first''; code:
\texttt{hpm\_better} \src{Scripts/Campaigns/2026-10-28_New_Campaign/code/core/v2_2_1/core/cpp/v3_util.h:120-126}, running best
updated on every fresh score event \src{Scripts/Campaigns/2026-10-28_New_Campaign/code/core/v2_2_1/core/cpp/v3_ledger.h:83-115},
table scores ingested first \src{Scripts/Campaigns/2026-10-28_New_Campaign/code/core/v2_2_1/core/cpp/v3_core.cpp:402-417}.
Lexicographic comparison is on sorted index lists with a proper prefix first
\src{Scripts/Campaigns/2026-10-28_New_Campaign/code/core/v2_2_1/core/cpp/v3_util.h:117-119}. The final local search starts from
the global best \src{Scripts/Campaigns/2026-10-28_New_Campaign/code/core/v2_2_1/core/cpp/v3_core.cpp:481}; within a scan the best
neighbour is chosen by $F$, exact ties by plain lexicographic order (not size first), and a move needs a strict $F$ increase
\src{Scripts/Campaigns/2026-10-28_New_Campaign/code/core/v2_2_1/core/cpp/v3_engine.h:377-425}. v3.1 adds an \emph{informational}
heuristic near-tie list at $\max(10^{-9},2\eta(\hat S))$ (e2) \src{NBH/branch_snapshots/MEET-v31/METHOD_v3.1_FINAL.md:181}.
Stores: top $40p+1$ by $F$ in HPM order; leader stores of $20p$ by $\log\gamma_0$ and by $\log\gamma_1$
\src{Scripts/Campaigns/2026-10-28_New_Campaign/code/core/v2_2_1/core/cpp/v3_ledger.h:73-81}.

\begin{stmt}{NBH/branch_snapshots/MEET-v31/METHOD_v3.1_FINAL.md:150}
\textbf{LS0: local search from the start mode.}
\begin{itemize}
\item ``After the final local search (and after the extra search from $c_4$ when $R'>4$), run one local search with the \S6.10 rules from
  the start mode $\tilde S_0=\arg\max\gamma_0$. $\tilde S_0$ is computed exactly by the block DP, with the size-then-lexicographic tie rule.''
\item ``LS0 reads no uniforms. Every model it scores joins the all-scored set and the stores.''
\item ``Models first scored by LS0 have provenance \texttt{ls0}. The flag \texttt{ls0\_improved} is set if LS0 raised the best score.''
\item ``\textbf{An LS0 improvement does not fire E3.} E3 refers only to the final local search from the HPM centre.''
\item ``An operational LS0 interruption means no complete output (GUARD-W, \#6). A completed LS0 capped at \texttt{local\_limit} still
  returns output (\#10).''
\item ``Start-mode ties: keep the enumeration fast path (cap $10^6$, \texttt{cpp/v3\_start.h:183}). On truncation, use a deterministic
  tie-aware DP that returns the prescribed smallest-size, lexicographically smallest maximiser without enumerating ties (e4, \#29).''
\end{itemize}
\end{stmt}
\#9: the bounded top store after LS0 is ``$\mathcal T_{\mathrm{new}}=\mathrm{top}_c(\mathcal T_{\mathrm{old}}\cup\mathcal S_{\mathrm{added}})$, not a
superset of $\mathcal T_{\mathrm{old}}$''; ``stores nested'' and ``bitwise fit'' are struck
\src{NBH/branch_snapshots/MEET-v31/METHOD_v3.1_FINAL.md:151}. \#10: \texttt{local\_certified} and card item L1 are evaluated on
the finally returned $\hat S$, from whichever search ended on it \src{NBH/branch_snapshots/MEET-v31/METHOD_v3.1_FINAL.md:152}.
LS0 cost: ``The runtime overhead is under 1 s per fit''; in the O-LS0 test export it changed the HPM in 0 of 2,394 fits
\src{NBH/branch_snapshots/O-LS0/REPORT.md:380} \src{NBH/branch_snapshots/MEET-v31/METHOD_v3.1_FINAL.md:48}. Rejected: an HPM--MPM
disagreement flag (``No provable power'') \src{NBH/branch_snapshots/MEET-v31/METHOD_v3.1_FINAL.md:277}.

\subsubsection{General-\texorpdfstring{$h$}{h} accumulator, \texorpdfstring{$\hat J$}{Jhat}, cross-block proxy}\label{sec:method-genh}
Definitions \rid{DF18}, \rid{DF19}, \rid{DF20}.
\begin{stmt}{NBH/branch_snapshots/MEET-v31/METHOD_v3.1_FINAL.md:159}
\textbf{Declaration.} ``𝓗 is a finite list of bounded model-space functions. The default list is fixed from $(X,y)$: model-size
indicators, and joint inclusion on neighbour-graph edges ($\le2pB$ pairs). A user may also declare $h$ measurable in $(X,y,\mathcal C)$,
fixed before stream 10. An $h$ defined from other calibration information (streams 0--9) is covered only if that information is
included in 𝓒 before it is frozen (enlarge 𝓒 when necessary); measurability in streams 0--9 need not imply measurability in the
compressed 𝓒. No final-run selection.'' \textbf{Accumulation.} ``$H_r(h)=(4N_e)^{-1}\sum_{q=1..4}\sum_ih(S_{rqi})$, taken after each complete
measured round ($\varphi=1$, equal weights after multinomial resampling). Report $\hat\mu(h)=\sum_r\hat Z_rH_r(h)/\sum_r\hat Z_r$ and the
equal-weight version.'' \textbf{Not allowed.} ``Products of marginal RB-PIPs are not joint estimates. v3.1 ships plain measured-round
particle averages for joint $h$ ($c_h=0$). Rao--Blackwellisation is used only for the marginal PIPs. No RB joint estimator ships.''
\textbf{Archive (optional).} ``The measured-round populations as a key → count multiset. An $h$ chosen adaptively after the run does not
inherit the fixed-$h$ guarantee.'' \textbf{Continuous parameters.} ``Use $\bar h(S)=E[h\mid S,X,y]$.'' \textbf{Descriptive only.} ``The
measured-round frequency of the returned HPM is \textbf{not} in 𝓗. It may be printed as a descriptive diagnostic with no theorem, next
to $e^{F(\hat S)}/\hat Z$.'' Placement in code: \texttt{cpp/v3\_engine.h:351--359}. The identity for general $h$ is a PROOF SKETCH (N1).
\end{stmt}
AMENDMENT\_HQ\_02 adds $d$ to the default 𝓗 (``raw d is recorded'', theorems use $(d-\min_\Omega d)/D$), the estimate
$\hat J=\sum_r\hat Z_rB_r/\sum_r\hat Z_r-(1/R')\sum_rA_r$ of the Jeffreys divergence $J=E_\pi d-E_{\pi_0}d$ with
$A_r=N^{-1}\sum_id(S^{(0)}_{ri})$ (initial start particles) and $B_r=(4N)^{-1}\sum_{k=1}^4\sum_id(S^{(k)}_{ri})$ (measured
endpoints), labelled ``Post hoc Monte Carlo estimate of Jeffreys divergence; not a certificate; floating-point error
unvalidated'' \src{NBH/branch_snapshots/MEET-v31/AMENDMENT_HQ_02_DEFINITIONS_from_MEET-THM.md:16-41}; and the cross-block proxy
$\hat\varepsilon_\times=\max_{b\ne c,\,j\in C_b,\,k\in C_c}|X_j^\top X_k|/(n-1)$ (\texttt{cross\_block\_rmax}) with companions
$\hat\kappa=\min_b\lambda_{\min}(G_{C_b}/(n-1))$, $\hat\delta=(K-1)\hat\varepsilon_\times/\hat\kappa$, $\hat{\bar u}=\sum_b\max_{U\subseteq C_b}x_b(U)$,
$\hat r=\hat\delta\hat{\bar u}/[(1-\hat\delta)(1-\hat{\bar u})]$, $\hat\eta=K\log(1/(1-\hat\delta))+\nu\log(1/(1-\hat r))$ on the stated domain, else
``n/a'' \src{NBH/branch_snapshots/MEET-v31/AMENDMENT_HQ_02_DEFINITIONS_from_MEET-THM.md:43-58}. ``Reporting fields 1--3 are specified,
not yet coded'' \src{NBH/branch_snapshots/MEET-v31/AMENDMENT_HQ_02.md:15}.

\subsubsection{Certificate card}\label{sec:method-card}
Definition \rid{DF17}; decision \rid{MD10}. The card is POST HOC, ``computed after the run and never feeds back into
calibration''. Abridged items below keep every quantifier; full text \src{NBH/branch_snapshots/MEET-v31/METHOD_v3.1_FINAL.md:212-255}.
\begin{stmt}{NBH/branch_snapshots/MEET-v31/METHOD_v3.1_FINAL.md:214-253}
\textbf{Header.} Every line is stated in exact arithmetic. Until SC($\eta_F$) is established by validated arithmetic and VAL-FP holds at
production size, every computed number carries the label ``numerical implementation of an exact-arithmetic certificate'', with $\eta$
printed. L1, L2 and R run under one declared certificate budget; on exhaustion the item reports ``unresolved''. A failed or unfinished
certificate never changes the fit's status. Under stacking there is one card per reported record.

\textbf{G1 (deterministic).} $\mathcal T$ is the deduplicated reported store and $Z_{\mathcal T}=\sum_{S\in\mathcal T}e^{F(S)}$. Then: $Z_{\mathcal T}\le Z$;
$\pi(S)\le e^{F(S)}/Z_{\mathcal T}$ for every $S$; $\pi(A)\le\sum_Ae^F/Z_{\mathcal T}$ for every $A\subseteq\mathcal T$; $\pi(S)/\pi(\hat S)\le e^{-g}$ on
$\mathcal T\setminus\{\hat S\}$, with $g=F(\hat S)-\max_{\mathcal T\setminus\{\hat S\}}F$; $\hat Z<Z_{\mathcal T}\Rightarrow\hat Z<Z$.

\textbf{G2 (one-sided, output-level).} Simultaneously for all $S$: $\pi(S)\le\min\{1,e^{F(S)}/\max(Z_{\mathcal T},0.05\hat Z)\}$;
$\pi(\Omega\setminus\mathcal T)\ge1-Z_{\mathcal T}/(0.05\hat Z)$. This is mass \textbf{outside the reported store} (never scored, or scored and
evicted). It is \textbf{not} a bound on never-scored mass. G2 fails only on $\{Z<0.05\hat Z\}$, and P(run returns output and G2 is false)
$\le0.05\cdot P(\mathrm{SMC}\cap\mathcal A)\le0.05$. This is a probability \textbf{joint with output, never conditional on completion}. [\dots]
If $Z_{\mathcal T}\ge0.05\hat Z$, G2 holds surely.

\textbf{L1 (local).} Evaluated on the finally returned $\hat S$ [\dots]. If \texttt{local\_certified}, then $F(S)\le F(\hat S)$ for every
feasible flip or swap. Report $g_{\mathrm{add}}$ ($+\infty$ if $|\hat S|=K$), $g_{\mathrm{del}}$ ($+\infty$ if $\hat S=\varnothing$),
$g_{\mathrm{swap}}$ ($+\infty$ if there are no swaps), $\rho_{\mathrm{loc}}=\pi(\hat S\mid\{\hat S\}\cup\mathcal N(\hat S))$.

\textbf{L2 (subsets).} $g_{\mathrm{sub}}=F(\hat S)-\max_{S\subsetneq\hat S}F(S)$ and $\sum_{S\subsetneq\hat S}\pi(S)/\pi(\hat S)$. With
$\max\varnothing=-\infty$: when $\hat S=\varnothing$, $g_{\mathrm{sub}}=+\infty$ and the relative subset mass is 0.

\textbf{L3 (best-of-scored).} If $\mathcal T=$ top-$(40p+1)$ of the realised scored set $\mathcal S$ and $|\mathcal T|\ge2$, then
$\pi(S)\le e^{-g}\pi(\hat S)$ for all $S\in\mathcal S\setminus\{\hat S\}$, with $g=F(\hat S)-F(\mathrm{top}[2])$. This holds under LS0 because
$\mathcal T_{\mathrm{new}}=\mathrm{top}_c(\mathcal T_{\mathrm{old}}\cup\mathcal S_{\mathrm{added}})$.

\textbf{R (optional).} A declared region rule $U(X,y,\hat S)\ni\hat S$ [\dots]. B\&B on $[\varnothing,U]\cap\Omega$, dropping leaves with
$|L|>K$, gives (a)--(c) as in PC2. ``Unresolved'' on budget exhaustion. In floating point, (a) needs $\tilde\Phi\le\tilde F(\hat S)-2\eta_F$.

\textbf{Card-level statement.} ``The probability that a run returns a card with a false G2, L or R claim is at most 5\%'', in exact
arithmetic. G1, L and R are true surely. [\dots] \textbf{Excluded:} C2; C5; the ``unscored mass'' wording; the B\&B cost remark; any
rigorous reading of \texttt{global\_certified}.
\end{stmt}
Cuts [\dots]: G2's premise list (PC1 claim 2 (i)--(iv), premise (i) = Lemma RO plus GUARD-W, a code-audit obligation);
L1's LS0 cap/interruption sentence (as \#10); R's printed node budget and $\tau$; the floating-point version of the card-level
statement (needs validated SC and VAL-FP). See the source span.

\subsubsection{Output convention, output event, unavailable outputs}
Definition \rid{DF21}; decision \rid{MD12}.
\begin{stmt}{NBH/branch_snapshots/MEET-v31/METHOD_v3.1_FINAL.md:168}
\textbf{Output convention.} ``Every estimate on the output event 𝒪 is a function of (X, y, seed, binary). The allocation and hardware
decide \textbf{delivery only} (stop-only guard, Lemma RO). Guarantees are stated for the potential (uninterrupted) run and as
\textbf{P\{𝒪, ·\}} (joint with output), never as P(· | 𝒪). P(no output) is an operational frequency that depends on the allocation and
hardware, and is reported with them (V13).''
\end{stmt}
\#15: ``On every return without complete output (\texttt{resource\_limit}, \texttt{schedule\_limit}, \texttt{numerical\_failure},
GUARD-W stops, LS0 stop, degenerate statuses), \texttt{fit\$log\_Z} and \textbf{every estimate field} ($\hat Z_r$, $\hat Z$, bounds, PIPs,
MPMs, $\hat\mu(h)$, card) are NA by exact name. [\dots] This fixes the partial-match hazard, where \texttt{fit\$log\_Z} silently
returned \texttt{log\_Zbar0}'' \src{NBH/branch_snapshots/MEET-v31/METHOD_v3.1_FINAL.md:162}. On \texttt{schedule\_limit} and
\texttt{numerical\_failure} the v3 rule is ``the best scored model is returned'', labelled incomplete
\src{Docs/Methodology/METHODOLOGY_v3.md:49-51}.

\subsection{The settings table and the \texorpdfstring{$(n,p)$}{(n,p)} rules}\label{sec:method-settings}

Definition \rid{DF25}; decision \rid{MD13}. Ground rule 9 requires every setting to be ``(a) modelling convention with a
stated reason; (b) explicit function of (n, p) or X; (c) declared resource budget; (d) numerical tolerance; (e)
arbitrary'', and ``\textbf{No setting is class (e) in the method's law}''
\src{NBH/branch_snapshots/MEET-v31/METHOD_v3.1_FINAL.md:287-294}; ``No separate p ≪ n / p ≫ n rule is needed: K, B, N₀, the
exact-branch rule and r\_n switch regime by formula'' \src{NBH/branch_snapshots/MEET-v31/METHOD_v3.1_FINAL.md:210}. The
constants paragraph: ``These values are conventions, not derived optima; changing them is not automatically
law-preserving.'' \src{NBH/branch_snapshots/MEET-v31/METHOD_v3.1_FINAL.md:208}. The table below transcribes all 49 rows of
METHOD \S5 (all six columns; markdown converted to LaTeX text) \src{NBH/branch_snapshots/MEET-v31/METHOD_v3.1_FINAL.md:296-346};
row $k$ is source line $297+k$, shown in parentheses in the first column. ``Numerical boundaries quoted from seats are
seat-derived unless marked [chair-checked].'' Code pointers inside the rows (e.g.\ \texttt{v3\_core.cpp:142}) are METHOD's own
and refer to the v2.1 tree (see the citation-version divergence in \S\ref{sec:method-impl}).

{\scriptsize
\setlength{\tabcolsep}{2.5pt}
\begin{xltabular}{\linewidth}{@{}>{\raggedright\arraybackslash}p{0.9cm}>{\raggedright\arraybackslash}X>{\raggedright\arraybackslash}p{1.25cm}>{\raggedright\arraybackslash}X>{\raggedright\arraybackslash}X>{\raggedright\arraybackslash}p{1.6cm}@{}}
\toprule
\# (line) & setting (value; source) & class & reason & behaviour as n, p grow & action\\
\midrule
\endhead
1 (298) & g = n, τ = 1 (M §5; \texttt{v3\_score.h:175}) & (a), (b) & unit information; ridge keeps m(S) finite for collinear X\textsubscript{S} & p-free; τ matters more for p ≫ n (collinear sub-models) & keep (Felipe) \\
2 (299) & IG(0.01, 0.01) (M §5) & (a) & proper, as the proofs require & pairwise effect ≤ n·b₀|1/Q₁ − 1/Q₂| on standardised data; relative to the gap it is not uniformly negligible (two equal-size, nearly collinear singleton models give a relative rate effect → 0.0142; review check, SOL1 \#43); negligibility needs residual-scale assumptions; not negligible at tiny n & keep; no universal "negligible" claim (R2\_C) \\
3 (300) & BB(1, p), w(k) = ½/C(2p−1, k), W\textsubscript{K} explicit (M §5) & (a), (b) & sparse, multiplicity-aware; the project theory is for it & defined for every finite p; dense growing truths are not its regime & keep \\
4 (301) & Cap K = min(p, ⌊n/2⌋) (M §3, §5) & (a), (b) & keeps near-interpolating models out & void for p ≤ n/2; binds for p ≳ n/2; K = 1 at n = 2, 3; s₀ > K unrecoverable by definition & keep \\
5 (302) & Input domain: n ≥ 2, p ≥ 1 (v3.1, proposed; current code rejects n < 3, \texttt{v3\_preprocess.R:103–107}) & (b) & standardisation and the target are defined at n = 2 & n = 2 → exact branch (K = B = 1, E = p + 1); n = 1 → \texttt{degenerate\_n}; n = 0 rejected & \textbf{e3} (§2 \#25) \\
6 (303) & Exact-constant column drop; \texttt{y\_constant}; \texttt{no\_predictors} (M §4) & (a) & removes undefined scales & regime-free & keep \\
7 (304) & Standardisation: max-abs, centre, SD (n−1) (M §4) & (a), (b) & affine and permutation invariance & regime-free & keep \\
8 (305) & Near-constant anchored two-pass route at √ε\textsubscript{mach} (\texttt{v3\_preprocess.R:47}) & (d) & cancellation protection, not screening & regime-free & keep; document (R2\_C) \\
9 (306) & Duplicate report |r| ≥ 1 − 10⁻¹⁰ (M §4 rule 5; \texttt{v3\_core.cpp:722}) & (d) & reporting only; duplicates are kept & regime-free & keep \\
10 (307) & B = min(⌊√n⌋, ⌊√p⌋, B*) (M §3) & (a), (b) & Felipe's rule; B ≤ K and B ≤ p & √p binds for p < n, √n for p > n; B = 1 at n = 2, 3 & keep \\
11 (308) & \textbf{B* = 16} (M §3, §8.7; \texttt{v3\_core.cpp:142}) & \textbf{(c)} & exhaustive-table budget: ≤ 4,096p scores and 65,536p bytes & binds once n, p ≥ 289; theorem SIZE condition: B = 16 covers groups 9..16 & \textbf{e1} (§2 \#23) \\
12 (309) & ℓ = ⌈log₂(n + p + 1)⌉ (M §3) & (b), (c) & local-search move cap and rebuild cap & 10–11 at the study, 20 at n + p ≈ 10⁶; not guaranteed enough for growing support (obligation) & keep \\
13 (310) & L = max(4, ⌈log₂(n + 1)⌉) (M §3) & (b), (c) & burn-in at φ = 1; floor binds for n ≤ 15 & grows with n, not with p; "L logarithmic versus t\textsubscript{mix} possibly polynomial" is V10 OPEN & keep \\
14 (311) & N₀ = max(2,500, n, p) (M §3) & (b), (c) & floor: the 2,500 is a declared budget. "Per-island RB-PIP MC SD ≤ 0.01" is an iid-sampling benchmark only (1/(2√N)); resampling creates dependence, so it is not a guarantee. Grows with dimension & floor binds at the study; N₀ = max(n, p) above 2,500 ⇒ Θ(p·max(n, p)) move steps per stage, Θ(p²) only when n = O(p); score-size costs extra (preflight reports it) & keep; C6 iff a particle-depth claim or any demonstrated fit with p > 2,500 (§3) \\
15 (312) & H₀ = ⌈512·N₀/2,500⌉ (M §3) & (b), (c) & declared calibration/final ratio 0.2048 & grows linearly; single-tier calibration/final = 0.0512 + 0.2048·I ≤ ≈ 2.1 at I = ℓ = 10 (seat-derived; ≤ ℓ rebuilds per tier) & keep; print the formula (§2 \#19) \\
16 (313) & T₀ = max(1,000, ℓ²) (M §3) & (b), (c) & finite stage ceiling; exceeding it ⇒ \texttt{schedule\_limit} & the ℓ² term binds only for n + p ≥ 2³¹; ideal-grid T* = O(log log p); as-coded T\textsubscript{n} OPEN (Conjecture P) & keep; \textbf{w4} (§2 \#33) \\
17 (314) & Tiers e ∈ \{0, 1\}, factor 4, ≤ 1 escalation (M §3, §6.9) & (c) & bounded escalation, decided before final randomness & ≤ ≈ 16× bridge work; no unbounded rescue; beyond → \texttt{schedule\_limit} & keep \\
18 (315) & R = 4 final islands; optional R′ = 4k at N\textsubscript{e}, declared before stream 10 (M §3; §2 \#7) & (c) & replication and parallel width & fixed 4 gives no growing-dimension accuracy; stacking at N\textsubscript{e} lowers the variance of the averaged evidence exactly ∝ 1/R′; for normalised PIPs and general h, 1/R′ is an asymptotic (delta-method) rate under moment conditions, not finite-R monotonicity & keep; \textbf{w7} \\
19 (316) & 4 diagnostic islands per rebuild (M §6.8) & (c) & E2 is an SD over 4 values & regime-free & keep \\
20 (317) & 4 measured rounds (M §6.11) & (c) & RB averaging after burn-in & constant; correlated rounds are not 4 replications & keep; general-h accumulator (§2 \#12) \\
21 (318) & Exact branch Σ\textsubscript{k≤K} C(p, k) ≤ (5H₀ + 4N₀)·p (M §6.2) & (b), (c) & enumerate when it costs about one SMC stage of all populations & fixed p is eventually exact as n grows (N₀ ≥ n); all p at n ≤ 5 [chair-checked]; n = 10: p ≤ 36; n ≥ 22: every p ≤ 17 is exact (a uniform sufficient range; the maximum for n ≤ 2,500, while larger n can admit more, e.g. n = 3,000, p = 18; review check, SOL1 \#33); growing p ≪ n need not be exact & keep; \textbf{w2} \\
22 (319) & Clustering: complete linkage on 1 − |r|, size-only cut, next-fit packing, lexicographic ties (M §6.3) & (b), (c) & deterministic X-dependent partition under the table budget & no correlation cut; O(p²) memory, O(p²)–O(p³) time; packing can merge separated groups & keep; \textbf{w5}; C11 (§3) \\
23 (320) & Tables (all subsets per block, exponent 1, full y) and start π₀ (M §6.3) & (a), (b) & exactly computable composite step, corrected at φ = 1 & constant B bounds each table; the number of blocks grows with p & keep; item 1 = Outcome 2: no depth or speed claim for the start \\
24 (321) & Neighbour graph: r\textsubscript{n} = √(2 log p / n), strict exceedance, forced best partner, ≤ B outgoing, symmetric union (M §6.4) & (b), (c) & chance-correlation scale of p columns; proposal heuristic & r\textsubscript{n} → 0 when log p = o(n); r\textsubscript{n} ≥ 1 iff p ≥ e\textsuperscript{n/2} (equality impossible for integer n, p) ⇒ 1-NN graph for p > 1; degrees may exceed B & keep; \textbf{w1}, \textbf{w3} \\
25 (322) & \texttt{cut\_group} diagnostic at ρ = r\textsubscript{n} (\texttt{v3\_core.cpp:760–770}) & (b) & descriptive only & can give large components or none & keep, descriptive \\
26 (323) & Move round m = p + B + 2, fresh random order; multinomial resampling each stage (M §6.5–6.6) & (b), (c) & reversible, self-adjoint kernel; fixed mutation allocation & Θ(p) sequential steps per particle per round: inherent to the kernel (I1) & keep; changing it is WORKHORSE \\
27 (324) & ESS threshold 0.90, pilot and grid (M §6.6–6.7; \texttt{v3\_calib.h:31, 79}) & \textbf{(c)}, with B's caveat & one-step dispersion budget (weight CV² ≤ 1/9). Caveat: "defines the law of 𝓒; carried by the theorems as a hypothesis (Theorem G needs the same threshold for pilot and grid)" & dimensionless; does not control accumulated error or omitted mass & keep; S3 R3-4 records (c) over B's (a) \\
28 (325) & 60 bisection halvings; endpoint snap 10⁻¹² (\texttt{v3\_calib.h:34–43, 82–94}) & (d) & 2⁻⁶⁰ below double resolution on [0, 1] & at very sharp crossings, progress may not resolve; failure is reported & keep \\
29 (326) & Pilot floor 1/T\textsubscript{e} (M §6.6) & (b), (c) & exploratory progress within the ceiling; forced steps counted and trigger escalation & shrinks as T\textsubscript{e} grows & keep \\
30 (327) & V = top 20p by γ₀ ∪ top 20p by γ₁ ∪ \{∅\}, |V| ≤ 40p + 1 (M §6.7) & (b), (c) & O(p) = size of a 1-flip neighbourhood; caps grid cost 60·|V|·T\textsubscript{e} & the retained fraction of model space vanishes; omitted mass not controlled & keep \\
31 (328) & Grid convergence = grid equality; ≤ ℓ rebuilds (M §6.7) & stopping rule (b), a deterministic rule (reclassified from the lead's (a), R2\_B row 20, by HQ amendment 01; values unchanged); cap (b)/(c) & a flat posterior keeps V changing by MC noise & convergence and P(S₀ ∈ V) OPEN (obligations) & keep; contract (§2 \#18) \\
32 (329) & E1 > 1 nat, E2 SD > 1 nat, E3, forced steps, grid failure/non-convergence, \texttt{local\_limit} (M §6.9) & (c) diagnostic criteria (S3 R3-4) & odds factor e; triggers define the law of 𝓒 & scale-free; no positive power as dimension grows & keep; labelled heuristics \\
33 (330) & Low-ESS flag: final stage < 0.5N (M §2) & (c), (b) & report flag & regime-free & keep \\
34 (331) & Top store 40p + 1; leader stores 20p (M §6.10) & (b), (c) & HPM survives eviction; near-ties, Z\_𝒯 and the card read it & linear in p; may truncate many near-equivalent models; restricted summaries labelled restricted & keep \\
35 (332) & Score cache 500p, clear when full (M §8.6) & (b), (c) & a hit returns identical bits; no effect on the law & key storage up to O(pK) & keep \\
36 (333) & Near-tie window 10⁻⁹ nat (M §6.10) & (d) & informational; the HPM uses exact equality & can fall below scorer error at large n; η depends on k, Q and conditioning, not on n alone, so a numerical crossing point needs a specified design and model (η ≈ 6·10⁻¹⁰ at n = 200; seat-derived) & keep + \textbf{e2} heuristic list (§2 \#24) \\
37 (334) & HPM tie rule: exact equality, then smallest |S|, then lexicographic (M §6.10) & (a) & argmax convention & scale-free & keep \\
38 (335) & MPM threshold ½; α = 0.05 for the Markov bound (M §6.11–6.12) & (a) & Hamming-loss rule; declared confidence & dimension-free; MPM may exceed K & keep \\
39 (336) & Q recompute below 10⁻⁸·yᵀy; Cholesky → refactor → augmented QR → \texttt{numerical\_failure} (M §6.13) & (d) & stable evaluation of one score & conditioning worsens with n, K; failure explicit & keep \\
40 (337) & Scorer error expression η (64ε…), used by \texttt{global\_certified} (\texttt{v3\_score.h:240}) & (d) as an indicator; (e) only if read as rigorous & the source calls it heuristic & no uniform validity established & \textbf{e5} (§2 \#27) \\
41 (338) & LS0 start-mode tie enumeration cap 10⁶ (\texttt{v3\_start.h:183}) & (d)/(c) fast path & DP argmax exact; only the tie choice can truncate & ties can proliferate with duplicates or symmetry & \textbf{e4} (§2 \#26) \\
42 (339) & Certificate budget (scores, B\&B nodes), region rule U, τ (card) & (c) declared certificate-region settings & PC2: Seat C's correction 6 upheld & "unresolved" on exhaustion; L2 needs 2\textsuperscript{|Ŝ|} − 1 scores & new in v3.1 (§2.1) \\
43 (340) & Default 𝓗: size indicators, neighbour-edge pairs ≤ 2pB (§2 \#12) & (b) & fixed from (X, y) via the graph & linear in pB & new in v3.1 \\
44 (341) & Allocation (default 8 GiB, 24 h), preflight, stop-only wall guard (M §2, §8.6) & \textbf{(c) declared input} & decides delivery only (Lemma RO) & a fixed allocation cannot support unrestricted growth (V13, by definition); the setup component alone, 8(3np + 3p²) bytes, gives p ≤ [−n + √(n² + 4·8·2³⁰/24)]/2 at 8 GiB (18,818 at n = 200; 11,398 at n = 20,000; review check, SOL1 \#39); this excludes all other allocations and workers and is not a feasible-fit threshold (seat-derived) & \textbf{e8}, \textbf{e6} (§2 \#22, \#28) \\
45 (342) & Seed ∈ [1, 2³¹ − 1]; stream map 0–9 (calibration), 10..9+R′ (final); 53-bit uniforms; integer dimensions & (c), (d) & reproducibility conventions; "bitwise" is per binary and BLAS & integer ranges can be exceeded at huge sizes & \textbf{e7} (§2 \#29); IND-T for 10..9+R′ (§6) \\
46 (343) & Workers: processes; worker count (§2 \#1–\#5) & (c) & delivery only; bitwise to serial & regime-free & new in v3.1 \\
47 (344) & Archive/export limits: V snapshot all if p ≤ 25, else ≤ 10⁴; R Gram copy only for p ≤ 2,000 (\texttt{v3\_calib.R:15–20}; \texttt{v3\_setup.R:9}) & (c) & storage choices; full V still builds the grid & truncation grows with p; never reconstruct a "full" V from a snapshot & keep; expose truncation \\
48 (345) & Other operational constants: exact-branch timing (200 scores, ≥ 64 columns), duplicate-check chunk 512, compiler-lock waits & (c) & profiling and workspace & hardware-dependent; outside the theory & keep \\
49 (346) & Gram-based scorer; O(p²) correlations and distances (M §8.1) & (c) implementation of the exact target & — & the dominant practical storage limit of v3.1; integer-dimension, stream and allocation ranges and checked-arithmetic rejection are also hard implementation limits (§2 \#29, row 45) & v3.2: C7, C11 (§3) \\
\bottomrule
\end{xltabular}
}
Per-regime summary (tiny $n$; fixed $p$, $n\to\infty$; growing $p\ll n$; $p\approx n$; $p\le n^A$; $p\le e^{Cn^\alpha}$; joint growth)
\src{NBH/branch_snapshots/MEET-v31/METHOD_v3.1_FINAL.md:350-377} and the adopted theory-scope sentence: ``The p ≤ n\textsuperscript{A}
scope of the selection theorems comes from the target theory (Felipe's prior and the proof technique), not from a sampler choice.
[\dots] Whether our prior admits it is OPEN, not impossible.'' \src{NBH/branch_snapshots/MEET-v31/METHOD_v3.1_FINAL.md:381}.
Method-induced restrictions: the cap excludes $s_0>K$; fixed $B$ limits whole-group recovery (as a single block) to $g\le B$;
bounded stores can omit posterior mass; logarithmic $\ell$ and $L$ may be short for some proof routes; $N_0\ge p$ and
full-coordinate rounds impose real computational growth \src{NBH/branch_snapshots/MEET-v31/METHOD_v3.1_FINAL.md:383-388}.

\subsection{The pipeline end to end}\label{sec:method-algo}

Assembled from M \S6.1 and \S7 \src{Docs/Methodology/METHODOLOGY_v3.md:133-149} \src{Docs/Methodology/METHODOLOGY_v3.md:293-306}
with the v3.1 additions \src{NBH/branch_snapshots/MEET-v31/METHOD_v3.1_FINAL.md:85}. Steps marked {\sffamily [v3.1]} are not in the
frozen code v2.2.1 (\S\ref{sec:method-impl}).

\begin{tcolorbox}[sopbase,colback=white,colframe=black!60,title={SMC-BVS v3.1$(X,y,s)$ --- pipeline}]
\small
\begin{enumerate}[leftmargin=*,label=\arabic*.]
\item \textbf{Checks and preprocessing} (\S\ref{sec:method-pre}); statuses \texttt{invalid\_input}, \texttt{no\_predictors},
  \texttt{y\_constant} {\sffamily [v3.1: \texttt{degenerate\_n}, \texttt{unsupported\_size}; $n=2$ admitted]}. Preflight $\to$ maybe
  \texttt{resource\_limit}.
\item \textbf{Budgets} $K,B,\ell,L,N_0,H_0,T_0$ from $(n,p)$. If $\sum_{k\le K}\binom pk\le(5H_0+4N_0)p$: \textbf{enumerate} (exact $Z$,
  PIPs, HPM in the strict order, near ties), status \texttt{exact}; report and stop.
\item \textbf{Setup}: Gram, correlations; complete-linkage blocks $C_1..C_J$ (size cut $B$, next-fit packing); exact tables
  $r_b(U)$ for all $U\subseteq C_b$ (all table scores enter the HPM ledger) {\sffamily [v3.1: tables in parallel]}; size profiles,
  $\bar Z_0$ by convolution, sampler of $\pi_0$; neighbour graph $N(\cdot)$.
\item \textbf{Calibration}, tier $e=0$ (streams 0 and 1--4): pilot ($H_e$ particles from $\pi_0$; adaptive ESS-0.90 schedule,
  floor $1/T_e$; multinomial resampling and one round per stage; $L$ burn-in + 4 measured rounds; local search) $\to$ $V$,
  preliminary grid $\to$ repeat $\le\ell$ times: 4 diagnostic islands ($H_e$ each) on the grid {\sffamily [v3.1: concurrent]}
  $\to$ local search $\to$ update $V$ $\to$ rebuild grid; stop when the grid repeats $\to$ triggers E1--E3, forced steps,
  grid failure/non-convergence, \texttt{local\_limit}.
\item If any trigger fired at $e=0$: re-preflight, then \textbf{one} escalation to tier 1 (streams 5 and 6--9), keeping score
  leaders. If grid construction fails at $e=1$: \texttt{schedule\_limit} (best scored model returned).
\item \textbf{Freeze} $\mathcal C$: grid, tier, $N_e$, $L$, inherited leaders/HPM {\sffamily [v3.1: and $R'$, declared before stream 10;
  default $R'=4$; optional $R'=4k$ or $R'_n=4\lceil\log_2(n+p+1)\rceil$]}.
\item \textbf{Final islands} $r=1..R'$ on streams $10..9+R'$ {\sffamily [v3.1: concurrent process workers, merged by Lemma S]}:
  $N_e$ fresh iid $\pi_0$ draws; for $t=1..T$: weights $e^{(\varphi_t-\varphi_{t-1})d}$, log-increment, multinomial resampling,
  one random-order round at $\varphi_t$; then $L$ burn-in and 4 measured rounds at $\varphi=1$ with RB recording of $o_j$
  {\sffamily [v3.1: $H_r(h)$ for $h\in\mathcal H$; $d$ means for $\hat J$]}.
\item \textbf{Final local search} from the global best (all scored models, all phases) {\sffamily [v3.1: scans in parallel; if
  $R'>4$ also from $c_4$]}; then {\sffamily [v3.1] \textbf{LS0}} from $\tilde S_0$.
\item \textbf{Report}: HPM $\hat S$ (strict order $F$, then $|S|$, then lex) with provenance; near ties; $\hat Z_r$, pooled $\hat Z$,
  $0.05\hat Z$; PIPs {\sffamily [v3.1: $\hat Z$-weighted primary, equal-weight secondary, both MPMs, $R_{\mathrm{eff}}$; $\hat\mu(h)$; $\hat J$,
  $\hat\varepsilon_\times$ and companions]}; flags; diagnostics {\sffamily [v3.1: POST HOC certificate card per record; on any return without
  complete output every estimate field NA]}.
\end{enumerate}
\end{tcolorbox}

\subsection{v3 \texorpdfstring{$\to$}{->} v3.1: what changed and what did not}\label{sec:method-diff}

``\textbf{Does the 2,400-fit study need a re-run? No.} The study stands as \textbf{v3.1 at R′ = 4, with POST HOC additions}''; its
saved sampler estimates are reusable with stated reporting limitations, and ``Full v3.1 records are not yet produced''
\src{NBH/branch_snapshots/MEET-v31/METHOD_v3.1_FINAL.md:45}. Unchanged (M \S\S): target (\S5); blocks, tables, composite start
(\S6.3); neighbour graph (\S6.4); move round (\S6.5); resampling (\S6.6); calibration law (\S6.6--6.9); budgets (\S3); HPM
decision (\S6.10); evidence estimator (\S6.12) \src{NBH/branch_snapshots/MEET-v31/METHOD_v3.1_FINAL.md:26-34}. The full
per-section map is \src{NBH/branch_snapshots/MEET-v31/METHOD_v3.1_FINAL.md:74-106}. Change classes: ADDITIVE (outputs on completed
fits unchanged or nested), POST HOC (computable from saved fits), LAW-PRESERVING (same law, different seed-to-uniform map),
WORKHORSE (changes the method's law) \src{NBH/branch_snapshots/MEET-v31/METHOD_v3.1_FINAL.md:16-20}. ``No change is WORKHORSE.''
\src{NBH/branch_snapshots/MEET-v31/MINUTES.md:28}. The table summarises METHOD \S2 rows (names shortened; see source lines).

{\footnotesize
\begin{xltabular}{\linewidth}{@{}l>{\raggedright\arraybackslash}X>{\raggedright\arraybackslash}p{2.3cm}>{\raggedright\arraybackslash}X@{}}
\toprule
\# & change & class & effect on the 2,400 fits (source line)\\
\midrule\endhead
0 & composite block start kept; item 1 = Outcome 2 (no depth/speed claim) & none & none \src{NBH/branch_snapshots/MEET-v31/METHOD_v3.1_FINAL.md:118}\\
1 & (a) concurrent final islands, Prop W & ADDITIVE (bitwise) & none; executed bitwise smoke + 3/3 \src{NBH/branch_snapshots/MEET-v31/METHOD_v3.1_FINAL.md:131}\\
2 & (b) concurrent diagnostic islands, Prop W-cal & ADDITIVE (bitwise) claimed from code reading only & none; W-cal not executed \src{NBH/branch_snapshots/MEET-v31/METHOD_v3.1_FINAL.md:132}\\
3 & (c) parallel local-search scans & ADDITIVE & none \src{NBH/branch_snapshots/MEET-v31/METHOD_v3.1_FINAL.md:133}\\
4 & (h) parallel block-table stage & ADDITIVE (bitwise) & none \src{NBH/branch_snapshots/MEET-v31/METHOD_v3.1_FINAL.md:134}\\
5 & process workers; counter rule & ADDITIVE (engineering) & none \src{NBH/branch_snapshots/MEET-v31/METHOD_v3.1_FINAL.md:135}\\
6 & GUARD-W semantics & ADDITIVE (spec) & none; the 6 \texttt{resource\_limit} fits stay stops \src{NBH/branch_snapshots/MEET-v31/METHOD_v3.1_FINAL.md:136}\\
7 & optional stacking $R'=4k$ at $N_e$; $c_4$ search; two records & ADDITIVE (nested) & none; study is $R'=4$ \src{NBH/branch_snapshots/MEET-v31/METHOD_v3.1_FINAL.md:142}\\
8 & LS0 from the start mode & ADDITIVE & 0 of 2,394 HPMs changed (test export); merged replay pending \src{NBH/branch_snapshots/MEET-v31/METHOD_v3.1_FINAL.md:150}\\
9 & D3c wording for LS0 (store = top$_c$, not superset) & spec text & none \src{NBH/branch_snapshots/MEET-v31/METHOD_v3.1_FINAL.md:151}\\
10 & certificate attaches to the returned $\hat S$ & ADDITIVE (spec) & none \src{NBH/branch_snapshots/MEET-v31/METHOD_v3.1_FINAL.md:152}\\
11 & $\hat Z$-weighted PIP primary; equal-weight secondary; both MPMs; $R_{\mathrm{eff}}$ & POST HOC & no refit; primary needs full rescoring (3 designs done: MPMs never differ, max$|\Delta\mathrm{PIP}|\le1.3\cdot10^{-3}$) \src{NBH/branch_snapshots/MEET-v31/METHOD_v3.1_FINAL.md:158}\\
12 & general-$h$ accumulator & ADDITIVE & not recoverable: no joint-summary claims \src{NBH/branch_snapshots/MEET-v31/METHOD_v3.1_FINAL.md:159}\\
13 & certificate card & POST HOC & computable where the top store was saved \src{NBH/branch_snapshots/MEET-v31/METHOD_v3.1_FINAL.md:160}\\
14 & $Z_{\mathcal T}$ and the underestimation certificate & POST HOC & as \#13 \src{NBH/branch_snapshots/MEET-v31/METHOD_v3.1_FINAL.md:161}\\
15 & unavailable output = explicit NA & ADDITIVE (interface) & only the 6 stopped fits' records change \src{NBH/branch_snapshots/MEET-v31/METHOD_v3.1_FINAL.md:162}\\
16--21 & output convention; card sentence; calibration contract; depth/cost; memory wording; \S9 limitation wordings & spec text & none \src{NBH/branch_snapshots/MEET-v31/METHOD_v3.1_FINAL.md:168-173}\\
22--29 & item-9 e8, e1, e2, e3, e4, e5, e6, e7 & spec text / POST HOC / ADDITIVE & none \src{NBH/branch_snapshots/MEET-v31/METHOD_v3.1_FINAL.md:179-186}\\
30--36 & wording fixes w1--w7 & spec text & none \src{NBH/branch_snapshots/MEET-v31/METHOD_v3.1_FINAL.md:192-198}\\
\bottomrule
\end{xltabular}
}

\begin{notproved}[what is not claimable for the study]
Quoting METHOD \S0 \src{NBH/branch_snapshots/MEET-v31/METHOD_v3.1_FINAL.md:52-61}: joint posterior summaries (``No joint-summary claims
for the 2,400 study fits''); the $\hat Z$-weighted (primary) PIPs and MPM on the full study until full rescoring (``So far only 3
designs have been rescored''; the equal-weight PIPs remain valid as the \textbf{secondary}); any stacking ($R'>4$) result; a
4-core speed-up (``The 1.9× figure is a 2-core measurement''); ``Bitwise'' for (b) until W-cal runs; ``Depth free of N, H, R'';
``LS0 as `run by the production core' until the merged-LS0 replay runs''; the certificate card wherever the top store was not
saved; ``Any depth or speed claim for the block start''.
\end{notproved}

Rejected changes (with reasons): backfit start (WORKHORSE, no exact sampler with known normaliser, no theorem); C8 split
$N_e$; C9 shorter/shared calibration; C10 parallel-in-particle kernel; HPM--MPM flag; truth-based gate for the primary PIP;
RB joint estimator; PSIS-$\hat k$/$\hat R$ flags; feeding certificates back \src{NBH/branch_snapshots/MEET-v31/METHOD_v3.1_FINAL.md:269-281}.
Engineering checks that gate claims (W-cal, merged-LS0 replay, GUARD-W audit, IND-T, VAL-FP, preflight, checked arithmetic,
full PIP rescoring, calibration timing) \src{NBH/branch_snapshots/MEET-v31/METHOD_v3.1_FINAL.md:392-407}.

\subsection{Implementation status (frozen code v2.2.1 vs spec)}\label{sec:method-impl}

All study fits used the frozen tree \texttt{Scripts/Campaigns/2026-10-28\_New\_Campaign/code/core/v2\_2\_1/}; ``v2.2.1 = v2.2 core,
byte-identical, plus one merged study layer'' and ``No numerical path changed''
\src{Scripts/Campaigns/2026-10-28_New_Campaign/code/core/v2_2_1/CHANGES_v2_2_1.md:3-25}; ``The v2.2.1 rerun reproduces the 2026-10-23
fits (\texttt{Results/SMC\_BVS\_v3\_production/}) exactly.'' \src{README.md:88}. Paths below are relative to that tree's
\texttt{core/}. ``Spec only'' = specified, no code in the repo; ``prototype'' = code on a \texttt{theory/*} branch (the
archive snapshots hold only .md/.tex/.lean).

{\footnotesize
\begin{xltabular}{\linewidth}{@{}>{\raggedright\arraybackslash}p{3.3cm}>{\raggedright\arraybackslash}X>{\raggedright\arraybackslash}p{3.1cm}@{}}
\toprule
component & where (code location, verified) & status\\
\midrule\endhead
preprocessing, statuses, budgets & \texttt{v3\_preprocess.R:19-33, 97-132} & implemented (v3 rules; $n<3$ rejected)\\
$n=2$, \texttt{degenerate\_n}, \texttt{unsupported\_size} (e3, e7) & --- & spec only\\
prior, score $m(S)$, $F(S)$ & \texttt{cpp/v3\_score.h:39-72, 175-238}; \texttt{cpp/v3\_engine.h:89-98} & implemented, matches\\
exact branch & \texttt{cpp/v3\_core.cpp:97-110, 526-570}; \texttt{v3\_exact.R} & implemented\\
blocks, cut, packing; neighbour graph & \texttt{cpp/v3\_blocks.h:39-173} & implemented, matches\\
tables, $\bar Z_0$, \texttt{draw\_pi0} & \texttt{cpp/v3\_start.h:73-176} & implemented (serial tables)\\
start mode $\tilde S_0$ (DP, cap $10^6$) & \texttt{cpp/v3\_start.h:178-240} (header only; production removed its Rcpp export) & present, unused in production; e4 tie-aware DP spec only\\
move kernels, random-order round & \texttt{cpp/v3\_engine.h:111-215} & implemented; N9 PASSED\\
resampling, island loop, RB recording & \texttt{cpp/v3\_engine.h:231-375} & implemented\\
pilot rule, $V$, grid, rebuild, escalation, freeze & \texttt{cpp/v3\_calib.h:25-98}; \texttt{v3\_calib.R:28-175} & implemented\\
final run, 4 islands, streams 10--13 & \texttt{cpp/v3\_core.cpp:421-520}; \texttt{v3\_run.R:84-90} & implemented, serial, $R=4$ fixed\\
(a) concurrent final islands & METH-AUDIT harness on branch \texttt{theory/METH-AUDIT} & prototype; executed bitwise 3/3\\
(b) concurrent diagnostic islands; (c) parallel scans; (h) parallel tables & --- & spec only (W-cal not executed)\\
stacking $R'$, $c_4$ search, two records & --- & spec only\\
LS0, \texttt{ls0\_improved}, provenance \texttt{ls0} & \texttt{ls0\_core.R}, \texttt{ls0\_exports.cpp} on branch \texttt{theory/O-LS0} & prototype (test export); not merged\\
HPM order, stores, near ties ($10^{-9}$) & \texttt{cpp/v3\_util.h:120-126}; \texttt{cpp/v3\_ledger.h:32-115}; \texttt{v3\_run.R:136-140} & implemented, matches\\
local search, \texttt{local\_certified} & \texttt{cpp/v3\_engine.h:377-425}; \texttt{v3\_run.R:141-147} & implemented\\
e2 heuristic near-tie list, $\eta(\hat S)$ & (\texttt{error\_bound} exists, \texttt{cpp/v3\_score.h:240-246}) & spec only\\
$\hat Z_r$, $\hat Z$, $0.05\hat Z$ & \texttt{cpp/v3\_engine.h:300-323}; \texttt{v3\_run.R:117-124} & implemented\\
equal-weight RB-PIP and MPM & \texttt{v3\_run.R:126-130} (named \texttt{pip\_primary}) & implemented (v3.1 \emph{secondary})\\
plain PIP & \texttt{cpp/v3\_engine.h:350-371}; \texttt{v3\_run.R:167-169} & implemented (v2.2 addition)\\
$\hat Z$-weighted PIP, $\mathrm{MPM}_{\hat Z}$, $R_{\mathrm{eff}}$ & \texttt{pip\_rescore.R} (MEET-v31 evidence, POST HOC) & POST HOC script; 3 designs rescored\\
general-$h$ (joint), $\hat J$, $\hat\varepsilon_\times$ & --- & spec only\\
$Z_{\mathcal T}$, certificate card & --- & spec only (POST HOC)\\
E1--E3, low-ESS flags; calibration triggers & \texttt{v3\_run.R:148-156}; \texttt{v3\_calib.R:58-118} & implemented\\
NA estimate fields & --- & spec only\\
stop-only runtime guard & \texttt{cpp/v3\_engine.h:336-339}; \texttt{v3\_fit.R:76-81} & implemented (serial core); v3.1 worker/LS0 audit open\\
preflight & \texttt{v3\_fit.R:16-46} & implemented: SMC memory only\\
\bottomrule
\end{xltabular}
}
Location sources: LS0 prototype \src{NBH/branch_snapshots/O-LS0/REPORT.md:69} \src{NBH/branch_snapshots/O-LS0/REPORT.md:95-97};
METH-AUDIT harness \src{NBH/04_code_and_files_index.md:366}; \texttt{pip\_rescore.R} \src{NBH/04_code_and_files_index.md:369};
production removals \src{Scripts/Campaigns/2026-10-28_New_Campaign/code/core/v2_2_1/core/README.md:15-19}.

\begin{divergence}[spec v3.1 vs code: primary PIP]
v3.1 makes the $\hat Z$-weighted PIP primary (\#11, \rid{MD08}). The frozen core computes only the equal-weight mean and calls it
\texttt{pip\_primary} (``equal-weight mean over islands (M §6.11)'')
\src{Scripts/Campaigns/2026-10-28_New_Campaign/code/core/v2_2_1/core/v3_run.R:129}; production removed ``secondary PIP computation and
outputs on both branches'' \src{Scripts/Campaigns/2026-10-28_New_Campaign/code/core/v2_2_1/core/README.md:16}, and the study adapter
stops if a secondary PIP is returned \src{Scripts/Campaigns/2026-10-28_New_Campaign/code/core/v2_2_1/study/adapter.R:229}. METH-AUDIT
flagged it: ``the Ẑ-weighted (secondary) PIP of METHODOLOGY §6.11 is not computed by the core''
\src{NBH/branch_snapshots/METH-AUDIT/REPORT.md:232}. Consequence: every study field or table named ``primary PIP'' is the v3.1
\emph{secondary} PIP. The v3.1 primary exists only via POST HOC rescoring, done on 3 designs
\src{NBH/branch_snapshots/MEET-v31/METHOD_v3.1_FINAL.md:54}.
\end{divergence}

\begin{divergence}[decision vs code: LS0 not in the production core]
\rid{MD03} decides ``LS0 in the shipped core''. In fact ``It is not yet in the production core, and the integrated replay is
pending'' \src{NBH/branch_snapshots/MEET-v31/METHOD_v3.1_FINAL.md:48}; METH-AUDIT: ``\textbf{LS0 is not in the production core}''
\src{NBH/branch_snapshots/METH-AUDIT/REPORT.md:235}. The O-LS0 check ran LS0 as a test export on the production core, and ``No
production SMC fits were rerun with LS0 integrated (\texttt{ls0\_enable\_integrated} exists in \texttt{ls0\_core.R} but was not
used)'' \src{NBH/branch_snapshots/O-LS0/REPORT.md:69}. The start-mode function is in the header but production removed ``the
start-mode Rcpp export'' \src{Scripts/Campaigns/2026-10-28_New_Campaign/code/core/v2_2_1/core/README.md:15}. Hence OP-LS0 is a
premise, not a code fact, and theorems ``about the shipped method'' that use LS0 (e.g.\ \rid{TH05} (b1-LH)) are about the
spec; merged-LS0 replay = engineering check 2 \src{NBH/branch_snapshots/MEET-v31/METHOD_v3.1_FINAL.md:399}.
\end{divergence}

\begin{divergence}[spec v3.1 vs code: islands, streams, parallel layer]
Code fixes \texttt{islands = 4L} \src{Scripts/Campaigns/2026-10-28_New_Campaign/code/core/v2_2_1/core/v3_run.R:88} and seeds final
islands on streams $9+r$, $r=1..4$ \src{Scripts/Campaigns/2026-10-28_New_Campaign/code/core/v2_2_1/core/cpp/v3_core.cpp:468-478},
serially; the stream map is 0--13 \src{Scripts/Campaigns/2026-10-28_New_Campaign/code/core/v2_2_1/core/v3_run.R:8-11}. No $R'$, no
$c_4$ search, no baseline/expanded records, no concurrency: ``The production core is serial throughout''
\src{NBH/branch_snapshots/METH-AUDIT/REPORT.md:172}. Prop W was executed with the METH-AUDIT harness, not the core; W-cal ``is
code reading only'' \src{NBH/branch_snapshots/METH-AUDIT/REPORT.md:216}.
\end{divergence}

\begin{divergence}[spec v3.1 vs code: unavailable outputs, statuses, guard]
\#15 requires every estimate field NA by exact name on any return without complete output. The core sets
\texttt{fit\$log\_Zbar0} after setup \src{Scripts/Campaigns/2026-10-28_New_Campaign/code/core/v2_2_1/core/v3_fit.R:146} and
writes \texttt{log\_Z} only on completion \src{Scripts/Campaigns/2026-10-28_New_Campaign/code/core/v2_2_1/core/v3_run.R:121}; on a
stop, R's partial matching makes \texttt{fit\$log\_Z} return \texttt{log\_Zbar0} --- the hazard \#15 names
\src{NBH/branch_snapshots/MEET-v31/METHOD_v3.1_FINAL.md:162}. Statuses \texttt{degenerate\_n} and \texttt{unsupported\_size} do not
exist; $n<3$ is \texttt{invalid\_input} \src{Scripts/Campaigns/2026-10-28_New_Campaign/code/core/v2_2_1/core/v3_preprocess.R:106}.
Also v3 \S2 defines \texttt{resource\_limit} as ``preflight predicts the memory or time allocation will be exceeded''
\src{Docs/Methodology/METHODOLOGY_v3.md:51-52}, while the code also returns it from a mid-run stop-only guard with partial,
labelled results \src{Scripts/Campaigns/2026-10-28_New_Campaign/code/core/v2_2_1/core/README.md:19}; the 6 study stops are such
wall-time stops \src{NBH/branch_snapshots/MEET-v31/METHOD_v3.1_FINAL.md:64}. v3.1's GUARD-W wording covers this (guard stop
$\Rightarrow$ no output).
\end{divergence}

\begin{divergence}[spec v3 vs code: preflight]
M \S8.6: ``estimate peak memory [\dots] and wall time [\dots]; if either exceeds the job allocation, return
\texttt{resource\_limit}'' \src{Docs/Methodology/METHODOLOGY_v3.md:321-324}. Code: ``Lead rulings [\dots] SMC preflight is memory
only'' \src{Scripts/Campaigns/2026-10-28_New_Campaign/code/core/v2_2_1/core/v3_fit.R:16}, with block count proxy
$\lceil p/B\rceil$ \src{Scripts/Campaigns/2026-10-28_New_Campaign/code/core/v2_2_1/core/v3_fit.R:28} --- the e6 issue (``use J ≤ p
blocks (not ⌈p/B⌉)'') \src{NBH/branch_snapshots/MEET-v31/METHOD_v3.1_FINAL.md:185}.
\end{divergence}

\begin{divergence}[reporting fields not coded]
General-$h$ for joint $h$ (sizes, neighbour-edge pairs, $d$), $\hat J$, $\hat\varepsilon_\times$ and companions, $Z_{\mathcal T}$ and the
certificate card, $R_{\mathrm{eff}}$, $\mathrm{MPM}_{\hat Z}$, the e2 heuristic near-tie list: no code in v2.2.1 (searched the tree);
``Reporting fields 1--3 are specified, not yet coded'' \src{NBH/branch_snapshots/MEET-v31/AMENDMENT_HQ_02.md:15}. The only
measured-round particle averages recorded are the marginal plain PIPs
\src{Scripts/Campaigns/2026-10-28_New_Campaign/code/core/v2_2_1/core/cpp/v3_engine.h:359-361}; ``The measured-round populations were
not saved'' \src{NBH/branch_snapshots/MEET-v31/METHOD_v3.1_FINAL.md:53}.
\end{divergence}

\begin{divergence}[citation versions: METHOD's code line numbers]
METHOD\_v3.1\_FINAL cites code lines that match the v2.1 tree \texttt{Scripts/SMC\_BVS/core/}, not v2.2.1: e.g.\ ``B* = 16 (M §3,
§8.7; \texttt{v3\_core.cpp:142})'' \src{NBH/branch_snapshots/MEET-v31/METHOD_v3.1_FINAL.md:308} is \texttt{B\_MAX} at line 142 of
\src{Scripts/SMC_BVS/core/cpp/v3_core.cpp:142} but at line 170 of
\src{Scripts/Campaigns/2026-10-28_New_Campaign/code/core/v2_2_1/core/cpp/v3_core.cpp:170}; \texttt{timings\$calibration} and
\texttt{timings\$final} are at \texttt{v3\_fit.R:124, 136} in v2.1 \src{Scripts/SMC_BVS/core/v3_fit.R:124} but at lines 158 and 170 in
v2.2.1 \src{Scripts/Campaigns/2026-10-28_New_Campaign/code/core/v2_2_1/core/v3_fit.R:158}. The header-file citations
(\texttt{v3\_score.h}, \texttt{v3\_start.h}, \texttt{v3\_calib.h}, \texttt{v3\_preprocess.R}) agree in both trees where checked. Not a
behavioural divergence; agents must resolve METHOD's \texttt{file:line} pointers against \texttt{Scripts/SMC\_BVS/core/}.
\end{divergence}

\textbf{Checked and found consistent} (no divergence): prior and score; cap $K$ in scorer, Gibbs step and local search;
block construction and tables; start convolution and sampler; neighbour graph; all four elementary steps and the random
order; multinomial resampling; pilot and grid rules; rebuild loop and single escalation; evidence estimator; HPM order and
running best over tables, calibration (inherited) and final phases; local-search rules; E1--E3 definitions.

\subsection{Method decisions MD01--MD13 (MEET-v31, 2026-11-03)}\label{sec:method-md}

Decision text = MINUTES \S2 ``Decision'' column (verbatim); reason = its ``Reason'' column (verbatim unless marked);
status = \texttt{ID\_REGISTRY} status column. Meeting: chair Opus 5.5 (non-voting), seats A, B (Fable 5.1), C (gpt-6-astra)
\src{NBH/branch_snapshots/MEET-v31/MINUTES.md:5-11}. Felipe's ground rules (DOSSIER \S1) are summarised in
\src{NBH/notation_and_conventions.md:11}; objectives: (1) accuracy in recovering the true posterior, (2) parallel
scalability, (3) BVS selection accuracy.

\newcommand{\mdbody}[4]{\ \textbf{#1}\quad \status{#2}\par
\noindent\emph{Decision.} #3\par\noindent\emph{Reason.} #4\par}

\par\smallskip\noindent\defid{MD01}\mdbody{MEET-v31 decision 1}{decided (MEET-v31 MINUTES §2); closed in: R1 (all three); R2 rule; R3-6 procedure; result §3}
{\textbf{Block (local-to-global) start kept.} The backfit start is rejected for paper 1. Arm D is v2 evidence only.
\src{NBH/branch_snapshots/MEET-v31/MINUTES.md:36}}
{``Ground rules 1, 5 and 7: no irrefutable case for dropping it; it carries the S1/S2 theorems. Backfit is WORKHORSE (full
re-run) with no theorem, and no exact sampler with a known normaliser was exhibited for it [\dots], so Ẑ unbiasedness is not
available. The earlier `headline gain capped at 6/2,400' rationale is withdrawn [\dots].'' Felipe: ``I'm open to dropping it but I
need a very serious, irrefutable argument'' \pos{1132}; ablation result Outcome 2 \src{NBH/03b_reasoning_and_traces.md:161-166}. METHOD
\S4 cites ground rules ``1, 5, 8''; HQ: ``both sets are correct in context''
\src{NBH/branch_snapshots/MEET-v31/AMENDMENT_HQ_01.md:141}.}

\par\smallskip\noindent\defid{MD02}\mdbody{MEET-v31 decision 2}{decided (MEET-v31 MINUTES §2); closed in: R1 (all three)}
{\textbf{HPM stays the only decision rule.} MPMs of both PIPs are reported as descriptions. No HPM-vs-MPM disagreement flag.
\src{NBH/branch_snapshots/MEET-v31/MINUTES.md:37}}
{``HPM is the 0--1 Bayes rule (T07), and the finite-n spine T38(0)/N200 is about it. A disagreement flag has no provable power
(ground rule 3).'' HQ: ``HPM stays primary, because it's the rule that is easiest to prove for an SMC output''
\src{NBH/03b_reasoning_and_traces.md:168-172}.}

\par\smallskip\noindent\defid{MD03}\mdbody{MEET-v31 decision 3}{decided (MEET-v31 MINUTES §2); closed in: R2 D3c}
{\textbf{LS0 in the shipped core,} after the final local search. Own provenance and \texttt{ls0\_improved} flag. Does not fire E3.
\src{NBH/branch_snapshots/MEET-v31/MINUTES.md:38}}
{``Makes T38(b1-LH) a theorem about the shipped method. ADDITIVE.'' Evidence: ``O-LS0: 0/2,394 HPM changes''. Earlier HQ advice:
``It closes the main theorem's biggest hole'', ``It costs almost nothing''; caveat ``most T38 premises fail at n=200 [\dots], so
LS0's value is asymptotic'' \src{NBH/03b_reasoning_and_traces.md:68-75}. See the LS0 divergence above.}

\par\smallskip\noindent\defid{MD04}\mdbody{MEET-v31 decision 4}{decided (MEET-v31 MINUTES §2); closed in: R1–R2 (D3e, D3f)}
{\textbf{Parallel layer} (each item ADDITIVE / bitwise unless stated): (a) concurrent final islands; (b) concurrent diagnostic
islands; (c) parallel scans with lexicographic reduction; (h) the block-table stage in parallel; process workers (threads need
private counters); GUARD-W ``any worker / LS0 / guard stop ⇒ no output'' (an LS0 stop is an operational interruption; a completed
LS0 capped at \texttt{local\_limit} still returns output, with L1 = n/a); (g) estimate fields NA on unavailable output.
\src{NBH/branch_snapshots/MEET-v31/MINUTES.md:39}}
{``Real issue: the parallel orientation (ground rule 2). The axis is islands and scans, not variables (ground rule 4).'' HQ:
``The real problem is in the code, not the method'' (production code serial everywhere)
\src{NBH/03b_reasoning_and_traces.md:153-159}.}

\par\smallskip\noindent\defid{MD05}\mdbody{MEET-v31 decision 5}{decided (MEET-v31 MINUTES §2); closed in: R2 D3a, D3b}
{\textbf{Stacking R′ = 4k at N\_e,} declared before stream 10 (streams 10..9+R′). N\_e is never split. Extra local search from the
R = 4 centre. Baseline and expanded records each carry their own flags and card. \src{NBH/branch_snapshots/MEET-v31/MINUTES.md:40}}
{``Keeps stacking ADDITIVE (nested).'' Evidence: ``PC1 claim 6: without the extra search the post-search HPM is \textbf{not}
monotone in R′ (COUNTEREXAMPLE). No listed theorem needs monotonicity; the extra search costs ≤ ℓ + 1 scans''.}

\par\smallskip\noindent\defid{MD06}\mdbody{MEET-v31 decision 6}{decided (MEET-v31 MINUTES §2); closed in: R2 D3d, D4; C6 trigger: HQ (AMENDMENT\_HQ\_01)}
{\textbf{Depth and cost are stated at the implemented (island) level.} ``Free of N, H, R'' is removed. Calibration is counted. C6
(per-particle streams) and C7 (Gram-free scorer) are the declared v3.2 path. C6 is implemented and law-checked iff paper 1 claims
particle-parallel depth \textbf{or} demonstrates any fit with p > 2,500. [\dots] Felipe decided (2026-11-03) that paper 1 claims
island-level depth only, and the study has p ≤ 1,000, so neither trigger currently fires.
\src{NBH/branch_snapshots/MEET-v31/MINUTES.md:41}}
{``Honest scaling claim.'' Evidence: calibration took 0.45×, 0.62× and 1.76× the serial final run on S1, S5 and S4. [\dots] =
the trigger's provenance (union of seat positions, HQ disposition), quoted in full in the registry row.}

\par\smallskip\noindent\defid{MD07}\mdbody{MEET-v31 decision 7}{decided (MEET-v31 MINUTES §2); closed in: R1–R2 (all three)}
{\textbf{Calibration kept unchanged, with a contract:} 𝓒 frozen before final randomness; validity holds for any frozen 𝓒; Theorem G
cap; ≤ ℓ rebuilds per tier (≤ 2ℓ with escalation); restricted ESS on V only; statuses distinguished. Shortening it (C9) is
rejected. \src{NBH/branch_snapshots/MEET-v31/MINUTES.md:42}}
{``Validity does not depend on calibration converging. The missing theorems concern efficiency and delivery (theorem
obligations).'' Shortening calibration ``weakens three existing theorems'' \src{NBH/03b_reasoning_and_traces.md:157}.}

\par\smallskip\noindent\defid{MD08}\mdbody{MEET-v31 decision 8}{decided (MEET-v31 MINUTES §2); closed in: R2 D5 (all three)}
{\textbf{Ẑ-weighted PIP is primary,} equal-weight is secondary, and both MPMs and R\_eff are reported. POST HOC.
\src{NBH/branch_snapshots/MEET-v31/MINUTES.md:43}}
{``The only PIP with a consistency-in-R theorem (Prop P(ii), T31(i)). Equal weighting keeps the single-island bias for every R. The
truth-based gate was withdrawn: closeness to the binary truth is objective 3, not objective 1.'' Evidence: 3-design rescoring
(MPMs never differ, max|ΔPIP| ≤ 1.3·10⁻³) \src{NBH/03b_reasoning_and_traces.md:174-176}.}

\par\smallskip\noindent\defid{MD09}\mdbody{MEET-v31 decision 9}{decided (MEET-v31 MINUTES §2); closed in: R1 (all three)}
{\textbf{Flags E1–E3, ESS and grid flags kept} as calibration triggers (they define the law of 𝓒) and as labelled heuristics. No
PSIS/R̂ flags. The certificate card never feeds back into calibration. \src{NBH/branch_snapshots/MEET-v31/MINUTES.md:44}}
{``Removing them is WORKHORSE. Their theory is negative (V12, E1c).''}

\par\smallskip\noindent\defid{MD10}\mdbody{MEET-v31 decision 10}{decided (MEET-v31 MINUTES §2); closed in: R2 D2; R3-1}
{\textbf{Certificate card (POST HOC):} PC2's minimal card. G1 [\dots]; G2 [\dots], joint with output, ≥ 95\%; L1–L3 [\dots]; R optional
(box). One declared certificate budget, with ``unresolved'' on exhaustion. L1 is evaluated on the returned Ŝ (after LS0). Each
number carries the label ``numerical implementation of an exact-arithmetic certificate'' until validated arithmetic (SC) and VAL-FP
exist; a long-double replay alone does not remove it. \src{NBH/branch_snapshots/MEET-v31/MINUTES.md:45} ([\dots] = item contents, given
in full in \S\ref{sec:method-card}.)}
{``The only diagnostics with a provable positive property (ground rule 3; HQ objection 8).''}

\par\smallskip\noindent\defid{MD11}\mdbody{MEET-v31 decision 11}{decided (MEET-v31 MINUTES §2); closed in: R2 D7; R3-3}
{\textbf{General-h estimator:} H\_r(h) after each complete measured round. Ẑ-weighted and equal-weight pools are reported. The default
𝓗 is model size and pairwise inclusion on neighbour-graph edges. h must be fixed from (X, y), or from (X, y, 𝓒) before stream 10.
Plain averages for joint h; Rao–Blackwellisation for marginal PIPs only. HPM terminal frequency is descriptive only. ADDITIVE.
\src{NBH/branch_snapshots/MEET-v31/MINUTES.md:46}}
{``Proposal §13.1 needs joint posterior summaries. The marginal RB product is not a joint estimate.'' The identity is a PROOF
SKETCH (obligation N1, \rid{OB01}).}

\par\smallskip\noindent\defid{MD12}\mdbody{MEET-v31 decision 12}{decided (MEET-v31 MINUTES §2); closed in: R2 item 9 e8}
{\textbf{Output convention:} every estimate on the output event 𝒪 is a function of (X, y, seed, binary). The allocation and hardware
decide only whether output is delivered (stop-only guard, Lemma RO). Guarantees are stated for the potential run and as P\{𝒪, ·\},
never conditional on completion. \src{NBH/branch_snapshots/MEET-v31/MINUTES.md:47}}
{``Ground rule 9 (\texttt{resource\_limit} vs `runnable from X, y and one seed').''}

\par\smallskip\noindent\defid{MD13}\mdbody{MEET-v31 decision 13}{decided (MEET-v31 MINUTES §2); closed in: R2–R3-4}
{\textbf{Item 9: generality in (n, p).} No setting is class (e) in the method's law. Fixes e1–e8 and wordings w1–w7 are adopted (list
in §4). No separate p ≪ n / p ≫ n rule is needed. \src{NBH/branch_snapshots/MEET-v31/MINUTES.md:48}}
{``Ground rule 9; METHODOLOGY\_v3 Requirement (4).'' Felipe's rule \pos{1214}: every setting a modelling convention with a reason, an
explicit function of $(n,p)$ or $X$, or a declared budget \src{NBH/03b_reasoning_and_traces.md:183-185}.}

\smallskip
Closure: ``\textbf{open method questions: none}'' after N9 passed \src{NBH/branch_snapshots/MEET-v31/MINUTES.md:150}. Binding-goal
rule: if an obligation tied to a binding proposal goal fails (N1 / general $h$, \S13.1; S4 weight stability, \S5.3; N9, \S5.1),
``the binding goal is not met by withdrawal: HQ escalates to Felipe for a method reconsideration or an explicit scope change''
\src{NBH/branch_snapshots/MEET-v31/MINUTES.md:161}. Growth rates in $(n,p)$ may be added to budgets, ``limited to growth rates only,
with no other method change'' (approved \pos{1304}) \src{NBH/notation_and_conventions.md:404}.

\begin{agentnote}[what a theorist must keep in mind about the method]
\begin{enumerate}[leftmargin=*]
\item \textbf{Object vs output.} The posterior $\pi$, $Z$ and the exact HPM/MPM are properties of the target (\S\ref{sec:method-target});
  the method outputs $\hat S$ (best of the \emph{scored} set $\mathcal S$), $\hat Z$, pooled PIPs. Statements about the method's
  output are about the \emph{potential (uninterrupted) run} and are transferred to delivery only as $P\{\mathcal O\cap\cdot\}$, never as
  $P(\cdot\mid\mathcal O)$ (\rid{MD12}, \rid{DF21}); Lemma RO ``does not imply $P(B\mid\mathcal O,\mathcal C)\le\alpha$'' and CEX-guard is a
  recorded COUNTEREXAMPLE \src{NBH/notation_and_conventions.md:433} (see \rid{TH15}).
\item \textbf{Spec vs code.} v3.1 items LS0, stacking, concurrency, $\hat Z$-weighted primary PIP, general $h$, card, NA fields are spec;
  the study's numbers come from v2.2.1 (\S\ref{sec:method-impl}). OP-LS0, GUARD-W, IND-T for $R'>4$ are \emph{premises}. Theorems
  ``about the shipped method'' with LS0 are about the spec until the merged-LS0 replay.
\item \textbf{Which ``primary''.} In code and study tables \texttt{pip\_primary} = equal-weight RB-PIP = v3.1 \emph{secondary}.
\item \textbf{Score.} $F(S)=\log\varpi(S)+\log m(S)$ includes the capped prior; HPM ties use $(F,\,-|S|,\,\text{lex})$; the local-search
  scan breaks exact ties by lex only. Notation $F$ clashes with V9's $F(S)$ and with a Frobenius $F_S$ (Appendix~\ref{app:notation}).
\item \textbf{Calibration is part of the law.} Triggers, ESS 0.90, grid equality and the inherited leaders define the law of $\mathcal C$;
  validity is conditional on \emph{any} frozen $\mathcal C$; convergence/delivery facts are OPEN obligations (\rid{OB05}).
\item \textbf{Randomness.} One serial stream per pilot/island; final streams $\ge10$ are fresh. IND-T is recorded as ``IND-T (PROVED,
  XF VERIFIED*, A-REV1 only)'', while ``IND-T for streams 10..9+R′: code audit, not done''
  \src{NBH/branch_snapshots/MEET-v31/THEOREM_REVISION_INPUTS.md:160} \src{NBH/branch_snapshots/MEET-v31/THEOREM_REVISION_INPUTS.md:231}.
  The independence step draws a fresh exact $\pi_0$ sample at every visit.
\item \textbf{Fixed budgets.} $R=4$, 4 measured rounds, floors 2,500/1,000, ESS 0.90, $\times4$ tier are declared (c) budgets; growth
  rates may be added only as explicit functions of $(n,p)$; anything else is a method change (WORKHORSE needs Felipe).
\item \textbf{Do not} claim depth/speed for the block start, depth free of $N,H,R$, joint summaries or primary PIPs for the 2,400
  study fits, or ``bitwise'' for concurrent calibration.
\end{enumerate}
\end{agentnote}
